% Hangboarding for Rock Climbing: A Literature Review
% Pablo Bustillo — April 2026
\documentclass[12pt,a4paper]{article}

\usepackage[T1]{fontenc}
\usepackage[utf8]{inputenc}
\usepackage{lmodern}
\usepackage{microtype}
\usepackage{geometry}
\geometry{margin=2cm}
\usepackage{setspace}
\singlespacing
\usepackage{parskip}
\setlength{\parskip}{4pt}
\usepackage{booktabs}
\usepackage{array}
\usepackage{longtable}
% \usepackage[backend=biber,style=numeric,sorting=none,maxbibnames=99]{biblatex}
% \addbibresource{bibliography.bib}
\usepackage{hyperref}
\hypersetup{colorlinks=true,linkcolor=blue,citecolor=blue,urlcolor=blue}
\usepackage{enumitem}
\usepackage{xcolor}
\usepackage{tcolorbox}
\tcbuselibrary{breakable}
\usepackage{caption}
\usepackage{amsmath}
\usepackage{amssymb}

% Evidence-tier boxes
\definecolor{evidencebg}{RGB}{230,245,230}
\definecolor{bioplausiblebg}{RGB}{255,250,220}
\definecolor{folklorecolor}{RGB}{240,230,250}
\definecolor{warncolor}{RGB}{255,230,230}
\definecolor{toverifybg}{RGB}{255,245,210}

% Box format: \begin{evidencebox}{Title · EVIDENCE_BACKED · ESS N}
\newtcolorbox{evidencebox}[2][]{colback=evidencebg,colframe=green!50!black,
  fonttitle=\bfseries,title={#2},#1,left=4pt,right=4pt,top=3pt,bottom=3pt}
\newtcolorbox{bioplausbox}[2][]{colback=bioplausiblebg,colframe=orange!70!black,
  fonttitle=\bfseries,title={#2},#1,left=4pt,right=4pt,top=3pt,bottom=3pt}
\newtcolorbox{folklorebox}[2][]{colback=folklorecolor,colframe=purple!60!black,
  fonttitle=\bfseries,title={#2},#1,left=4pt,right=4pt,top=3pt,bottom=3pt}
\newtcolorbox{toverifybox}[2][]{colback=toverifybg,colframe=orange!80!red,
  fonttitle=\bfseries\color{orange!80!red},title={#2},#1,left=4pt,right=4pt,top=3pt,bottom=3pt}
\newtcolorbox{warningbox}[2][]{colback=warncolor,colframe=red!60!black,
  fonttitle=\bfseries\color{red!70!black},title={#2},#1,left=4pt,right=4pt,top=3pt,bottom=3pt}
\newtcolorbox{pendingbox}[1][]{colback=warncolor,colframe=red!60!black,
  fonttitle=\bfseries\color{red!70!black},title=Content Pending LLM Synthesis,#1,
  left=4pt,right=4pt,top=3pt,bottom=3pt}

% Status
\newcommand{\toverify}{\textbf{To Verify}}
\newcommand{\bioplausible}{\textbf{Bioplausible}}
\newcommand{\folkloreverified}{\textbf{Folklore Verified}}
\newcommand{\evidencebacked}{\textbf{Evidence-Backed}}

\title{\textbf{Training for Rock Climbing}\\[0.3em]
\large A Literature Review of Scientific Evidence and Practitioner Consensus}
\author{Pablo Bustillo}
\date{April 2026}

\begin{document}
\maketitle
\tableofcontents
\newpage

%---------------------------------------------------------------------
\begin{abstract}
This manuscript reviews the scientific evidence and practitioner consensus relevant to rock climbing training across five domains: hangboarding, weekly training integration, supplementation, injury management, and recovery. Evidence is graded using an Evidence Strength Score (ESS 0--10) where ESS 10 = well-powered climbing RCT with functional outcomes and ESS 1--2 = anecdotal or mechanistic-only support; every actionable recommendation carries a status label (EVIDENCE\_BACKED, BIOPLAUSIBLE, FOLKLORE\_VERIFIED, or TO\_VERIFY).

For \textbf{hangboarding}, ten intervention trials (nine RCTs and one controlled study, combined $n \approx 230$), one retrospective cohort ($n = 527$), and one systematic review and meta-analysis \mbox{(Stien et~al.\ 2023, $n = 110$)} support three principal conclusions: (1) supplemental hangboarding 2--3$\times$/week consistently improves maximal finger strength (MFS) and grip endurance within 4--10 weeks in intermediate-to-advanced male climbers (meta-analytic SDM~= 0.41 for MFS, 95\%~CI 0.03--0.80; SDM~= 0.57 overall, 95\%~CI 0.24--0.91); (2) $\geq$80\% MFS protocols maximise MFS, 60--80\% MFS repeater protocols maximise stamina; (3) daily low-intensity hangs produced a positive but non-significant trend versus climbing-only controls ($+2.5\%$, $p = 0.12$). No trial has demonstrated a statistically significant climbing-grade improvement attributable to hangboarding alone; all effect sizes are on surrogate dynamometric measures.

For \textbf{weekly integration}, session sequencing rules are predominantly BIOPLAUSIBLE (ESS 2); the $\geq$48\,h inter-session minimum is a study-design convention, not an empirically validated prescription. For \textbf{supplementation}, no supplement has been tested with a functional climbing primary outcome; collagen + vitamin C and beta-alanine carry the strongest mechanistic and surrogate evidence for connective-tissue support and endurance tasks, respectively (ESS 3--7); caffeine showed null results in the only climbing-specific RCT. For \textbf{injury management}, A2 pulley conservative rehabilitation with pulley-protection splint (PPS) is the evidence-supported first-line approach (ESS 5; 88.4\% return to prior level at mean 8.8\,months in a prospective series); Grade I--III injuries are managed conservatively; surgery is reserved for Grade IVa/IVb or FLIP complications. The section additionally covers flexor tenosynovitis (second most common finger injury; favourable natural history, ESS 5), PIP capsulitis/synovitis (third most common; 72\% return with staged algorithm, ESS 5), lumbrical tears (climbing-specific; lumbrical stress test pathognomonic, ESS 5), adolescent physeal injuries (most common injury under age 16--17; pulley framework contraindicated), and NSAIDs (mechanistic caution for prolonged use; no climbing-specific trial). Experience-dependent risk is observationally established (ESS 5). For \textbf{recovery}, sleep $\geq$8\,h and strategic cold water immersion carry the strongest evidence base across sports, with no climbing-specific RCT for any modality.

The literature is uniformly limited by small samples (median PEDro 5--7/10), male-advanced cohorts, short follow-up ($\leq$10 weeks), and surrogate outcomes. Female-cohort data, long-duration trials with injury tracking, and any trial measuring grade improvement as a primary outcome are absent across all five domains.
\end{abstract}

%---------------------------------------------------------------------
\section{Introduction}

Hangboarding involves isometric dead-hangs from edges, pockets, or slopers to overload the finger flexors beyond what routine climbing provides. MFS and forearm endurance are among the strongest measurable predictors of climbing performance in trained cohorts: Pearson correlations of $r = 0.70$--$0.85$ between MFS (or grip endurance) and redpoint grade are consistently reported \cite{Wat04, MSB+07, FSS+16}. General upper-body resistance training does not produce equivalent finger-flexor adaptations, establishing hangboarding as a distinct training stimulus. The first hangboard RCTs appeared circa 2012 \cite{LRGB12, MKL15}; the current literature comprises ten intervention trials and one large retrospective cohort \cite{LRGB12, MKL15, LRGB19, LL19, MSA+21, DO22, HO22, DO24, LAS24, RLPCBS25, GKAO24}, with intervention $n = 11$--54 and cohort $n = 527$, durations of 4--10 weeks. A systematic review and meta-analysis by Stien et al.\ \cite{SO23} synthesised five of these trials ($n = 110$). Universal limitations include lack of allocation concealment, male-advanced-only samples, and outcomes measured on the hangboard rather than on-wall performance.

%---------------------------------------------------------------------
\section{Methods}

\subsection*{Search Strategy and Scope}

This narrative review searched PubMed and Google Scholar across five domains: (1)~\textbf{hangboard and finger-strength intervention studies} (primary literature base); (2)~\textbf{connective tissue supplementation} (collagen peptides, gelatin, omega-3, creatine, $\beta$-alanine, caffeine, magnesium); (3)~\textbf{finger and upper-limb injury management} (pulley injuries, flexor tendon strains, skin trauma); (4)~\textbf{athletic recovery modalities} (sleep, HRV, contrast water therapy, soft-tissue work, nutrition timing); (5)~\textbf{weekly training integration and periodization} (concurrent training, session sequencing, deload protocols). Searches were conducted to April 2026. Systematic reviews, meta-analyses, and practice guidelines were additionally sourced from Cochrane, GRADE working group, and Physiotherapy Evidence Database (PEDro). Practitioner claims from widely-used coaching resources (\textit{Training for Climbing}, \textit{Rock Climber's Training Manual}, Lattice Training, Dave MacLeod, Crimpd, \texttt{r/climbharder}) are reported and classified by evidence tier.

\subsection*{Inclusion and Exclusion Criteria}

\textbf{Peer-reviewed studies} were included if they: (i) used a hangboard, no-hang device, or climbing-specific resistance protocol as the primary intervention with $\geq$1 finger-strength, endurance, or tendon-structural outcome, or (ii) examined a supplementation, recovery, or rehabilitation modality with direct relevance to a climbing-specific tissue (finger flexors, annular pulleys, skin) or performance outcome. Case series and observational studies were included when RCT-level evidence was absent (injury epidemiology, diagnostic accuracy). Foundational biomechanical and physiological studies were included to contextualise training and rehabilitation rationale.

\textbf{Adjacent-sport transfer evidence} (studies in general resistance-training, running, cycling, or team-sport populations) was included only when: (a) no climbing-specific controlled study existed for the question; (b) the tissue, loading modality, or physiological mechanism was explicitly comparable; and (c) the transfer assumption was stated and rated using ESS indirectness penalties (typically ESS 3--4 vs.\ ESS 7--8 for direct climbing evidence). All claims derived from adjacent-sport evidence carry a transfer-uncertainty qualifier.

\textbf{Folklore and practitioner claims} were included when they originated from $\geq$3 independent climbing-community sources (see Claim Status Labels below). Sources with $<$3 independent attestations are classified \toverify and excluded from settled guidance.

\subsection*{Claim Status Labels}

Each claim carries one of four labels:
\begin{itemize}[noitemsep]
  \item \textbf{\evidencebacked}: peer-reviewed intervention evidence with transparent design and explicit bias context.
  \item \textbf{\bioplausible}: mechanistically plausible, indirect or incomplete intervention support; mechanism stated; direct RCT gap noted.
  \item \textbf{\folkloreverified}: community-held belief with prevalence verified from $\geq$3 independent climbing-community sources (e.g., named coaches, published books, subreddit wiki).
  \item \textbf{\toverify}: provenance unclear or fewer than 3 independent source confirmations; must not be treated as settled guidance.
\end{itemize}

\subsection*{Evidence Strength Score (ESS)}

Each actionable recommendation carries an ESS from 0--10 reflecting study design, sample relevance to climbers, risk of bias, outcome directness, and transfer indirectness:

\begin{center}
\small
\begin{tabular}{>{\raggedright}p{1.2cm} >{\raggedright\arraybackslash}p{12cm}}
\toprule
\textbf{ESS} & \textbf{Criteria} \\
\midrule
10 & Well-powered climbing RCT, low risk of bias, functional outcome (redpoint grade, route completion, injury incidence) \\
9  & Strong climbing RCT with minor limitations \\
8  & Climbing RCT with moderate limitations \\
7  & Controlled non-randomised climbing intervention or high-quality climbing systematic review / meta-analysis \\
6  & Prospective climbing cohort or repeated-measures climbing intervention with acceptable controls \\
5  & Climbing observational evidence with meaningful confounding \\
4  & High-quality adjacent-sport RCT with explicit transfer uncertainty (e.g., Achilles tendinopathy $\to$ A2 pulley) \\
3  & Observational adjacent-sport evidence with substantial indirectness \\
2  & Mechanistic rationale, expert interpretation, or in-vitro evidence; no direct intervention support \\
1  & Anecdotal, single-source, or highly limited evidence \\
0  & Unreliable, unresolvable, or critically biased source basis \\
\bottomrule
\end{tabular}
\end{center}

\noindent ESS is reduced by one level when: (i) primary outcome is a surrogate biomarker (e.g., P1NP, serum markers) rather than a functional or structural endpoint; or (ii) a study is single-source and has not been independently replicated. These reductions are documented at each recommendation site.

\subsection*{Endpoint Classification}

All recommendation sites explicitly distinguish \textbf{surrogate} outcomes (biomarkers, dynamometric measures, VAS scores, in-vitro proxies) from \textbf{functional} outcomes (on-wall climbing performance, injury incidence, return-to-sport, grade improvement). Surrogate outcomes are labelled \textit{(surrogate)} at first use; functional outcomes are labelled \textit{(functional)}. Recommendations derived solely from surrogate-endpoint studies are assigned reduced recommendation confidence regardless of ESS.

\subsection*{Synthesis Process}

Initial multi-domain literature queries used eight free-tier frontier LLMs to maximise recall across specialist topics (supplements, injury, recovery, weekly programming). Synthesis, claim classification, and ESS assignment were performed by the author with AI-assisted orchestration, with all claims reconciled against primary sources. Single-source claims carry an explicit footnote and are excluded from any recommendation presented as settled guidance. Named model brands do not appear in manuscript body text; methods may state that free-tier frontier LLMs were used for feedback and that synthesis was programmatically orchestrated. See Appendix~\ref{app:toverify} for the consolidated \toverify tracker.


\section{Hangboarding}
\label{sec:hangboarding}
%---------------------------------------------------------------------
\subsection{Biomechanical and Physiological Foundations}

In a full-crimp grip (DIP joint hyperextended), biomechanical modelling estimates A2 pulley forces at approximately $2.7\times$ the external fingertip force (${\sim}255$~N at ${\sim}96$~N fingertip load) \cite{VQLVM06}, while A2 loading in crimp is approximately 36 times greater than in the open-hand (slope) position---a grip-type comparison, not a load multiplier. Cadaveric studies confirm high A2 pulley loads in crimp (380--700~N across finger positions) \cite{Sch01}, accounting for the predominance of A2 ruptures in climbing. Half-crimp and open-hand grips generate substantially lower pulley loads, underpinning standard advice against routine full-crimp hangboarding. Middle and ring fingers bear disproportionately higher loads than index and little fingers \cite{VQPC08}$^{\dagger}$, consistent with clinical injury patterns.

Climbing physiology reviews \cite{Wat04, MSB+07, FSS+16} establish three points central to hangboard training design: MFS is the single strongest measurable predictor of redpoint grade in trained cohorts; forearm endurance and oxidative capacity independently predict sustained-route performance; and these qualities are specific enough that general upper-body resistance training does not substitute for direct finger loading.

\smallskip
\noindent$^{\dagger}$\textit{Vigouroux et al. \cite{VQPC08} is a conference abstract; the force-distribution finding is consistent with the broader biomechanical literature but should be treated as supporting rather than primary evidence.}

%---------------------------------------------------------------------
\subsection{Intervention Studies}

\subsubsection{Study Summary}

Table~\ref{tab:studies} summarises all primary intervention studies.

\begin{longtable}{>{\raggedright\arraybackslash}p{2.6cm}
                  >{\raggedright\arraybackslash}p{1.9cm}
                  >{\raggedright\arraybackslash}p{3.4cm}
                  >{\raggedright\arraybackslash}p{4.0cm}
                  >{\raggedright\arraybackslash}p{1.8cm}}
\caption{Hangboard intervention studies, 2012--2025.}\label{tab:studies}\\
\toprule
\textbf{Study} & \textbf{Design/$n$} & \textbf{Protocol} & \textbf{Main finding} & \textbf{Cite} \\
\midrule
\endfirsthead
\multicolumn{5}{c}{\textit{(continued)}}\\
\toprule
\textbf{Study} & \textbf{Design/$n$} & \textbf{Protocol} & \textbf{Main finding} & \textbf{Cite} \\
\midrule
\endhead
\bottomrule
\endfoot
\bottomrule
\endlastfoot

López-Rivera \& G-B 2012 & 2-arm RCT; $n=16$ elite & 8~wk, 2$\times$/wk: max weight on 18~mm vs.\ minimal edge, same duration & Max-weight arm superior on strength; minimal-edge arm superior on endurance transfer & \cite{LRGB12} \\[3pt]

Medernach et al.\ 2015 & 2-arm RCT; $n=23$ adv.\ male & 4~wk, 3$\times$/wk fingerboard vs.\ bouldering only & FB: $+2.5$~kg grip, $+5$--7~s DH, $+26$~s IFH vs.\ control & \cite{MKL15} \\[3pt]

López-Rivera \& G-B 2019 & 3-arm RCT; $n=26$ advanced & 8~wk, 2$\times$/wk: MaxHangs / IntHangs / Combo & Endurance: IntHangs $+45\%$ (ES~1.0), MaxHangs $+34\%$ (ES~0.6), Combo $+7\%$; MFS: MaxHangs $+28\%$ only & \cite{LRGB19} \\[3pt]

Levernier \& Laffaye 2019 & Controlled; $n=23$ elite & 4~wk, 2$\times$/wk unilateral hangs (half-crimp $+$ slope) & $+8\%$ MFS; $\uparrow$RFD; control unchanged & \cite{LL19} \\[3pt]

Mundry et al.\ 2021 & 3-arm RCT; $n=30$ int--adv & 8~wk, 2--3$\times$/wk: added weight on 20~mm / progressive edge reduction / control & Added weight $>$ control ($p=0.032$, ES~$0.36$); edge reduction $=$ control & \cite{MSA+21} \\[3pt]

Devise et al.\ 2022 & 4-arm RCT; $n=54$ experienced & 4~wk, 2$\times$/wk: F100 / F80 / F60 / control; 12~mm instrumented hold & F100$+$F80: MFS $+12$--14\%; F80 stamina score: 29.5\%$\to$60.4\% (+31~pp); F60: 32.6\%$\to$53.5\% (+21~pp); F80 improves all three outcomes & \cite{DO22} \\[3pt]

Hermans et al.\ 2022 & 2-arm RCT; $n=35$ int--adv & 10~wk, 2$\times$48~min/wk progressive repeaters (BM1000, 23~mm) vs.\ climbing only & Peak force $+17\%$, avg force $+18\%$ (within-group); RFD $+28\%$ within HBT group but between-group difference not significant ($p=0.056$); DH $+12\%$ within HBT & \cite{HO22} \\[3pt]

Dindorf et al.\ 2024 & 2-arm RCT; $n=18$ completers (14M/4F) int--adv & 7~wk: NMES-assisted fingerboard 2$\times$/wk vs.\ conventional fingerboard 3$\times$/wk & NMES non-inferior to conventional training for finger flexor endurance despite ${\sim}33\%$ lower volume; no significant between-group difference & \cite{DO24} \\[3pt]

Gilmore et al.\ 2024 & Retrospective cohort; $n=527$ Crimpd users & Abrahangs ($\geq3\times$/wk, ${\sim}40\%$ MFS) vs.\ Max Hangs ($\geq0.5\times$/wk, 85--95\%) vs.\ Both vs.\ Climbing Only & Abrahangs $+2.5\%$, Max Hangs $+3.2\%$, Both $+5.8\%$, Climbing Only $+0\%$ strength-to-weight; self-selected, self-reported loads & \cite{GKAO24} \\[3pt]

Langer et al.\ 2024 & RCT; $n=31$ female ($n=26$ completed) & 5~wk: off-wall (fingerboard, BM1000, progressive levels) vs.\ on-wall vs.\ control & Trends toward fingerboard superiority on strength (BF$_{10}=1.16$, anecdotal); on-wall preferred for enjoyment; no significant between-group strength differences & \cite{LAS24} \\[3pt]

Ruiz-López et al.\ 2025 & RCT; $n=11$ experienced (of 20 enrolled) & 4~wk: weighted dead-hang (HIMA, 85\%, 14~mm) vs.\ active-pull dynamometer (PIMA) & Both groups significant gains in peak force, RFD$_{200}$, and max hang time ($p \leq 0.05$); no between-group differences; PIMA reduced hand asymmetry & \cite{RLPCBS25} \\[3pt]

\end{longtable}

\subsubsection{Max Hangs vs.\ Intermittent Hangs (Repeaters)}
\label{sec:maxvsrep}

\begin{evidencebox}{MaxHangs vs.\ IntHangs: Intensity-Specificity (López-Rivera 2019) · \evidencebacked · ESS 7}
López-Rivera \& González-Badillo \cite{LRGB19} randomised 26 advanced climbers (${\sim}7c{+}/8a$, mean age 31.7~yr) to MaxHangs (maximal added weight on 18~mm then minimal edge), IntHangs (4--5 sets $\times$ 10~s on / 5~s off on minimal edge), or Combination (4 weeks each in sequence). IntHangs produced the largest endurance gain ($+45\%$, ES~1.0); MaxHangs the largest MFS gain ($+28\%$, ES~0.6); the Combination group improved endurance by only $+7\%$ and gained no additional MFS. Max hangs maximise MFS; repeaters maximise endurance; sequential combination within one mesocycle underperforms both. \textit{Endpoint: dynamometric surrogate (time-to-failure dead hang, grip endurance) — no on-wall performance or climbing grade outcome measured.}
\end{evidencebox}

\begin{bioplausbox}{Concurrent Training Interference: Mechanism for Poor Combination Result · \bioplausible · ESS 2}
The inferior Combination result is consistent with strength--endurance concurrent-training interference but may equally reflect insufficient per-protocol volume when both are compressed into one mesocycle. This mechanism has not been tested directly in climbing studies.
\end{bioplausbox}

Devise et al.\ \cite{DO22} found that 80\% MFS protocols simultaneously improved MFS ($+12$--14\%), stamina (score rising from 29.5\% to 60.4\%, a gain of $+31$ percentage points), and endurance, making high-submaximal loading the pragmatic default for athletes training both qualities in parallel. The 60\% MFS protocol improved stamina by $+21$ percentage points and produced smaller MFS gains, consistent with an intensity--adaptation specificity relationship.

\subsubsection{Added Weight vs.\ Reduced Edge Depth}

\begin{evidencebox}{Added Weight Outperforms Progressive Edge Reduction (Mundry 2021) · \evidencebacked · ESS 7}
Mundry et al.\ \cite{MSA+21} found that progressive added weight on a 20~mm edge significantly outperformed both edge reduction and control on overall grip strength ($p=0.032$, ES~$0.36$, significant on 3 of 7 pinch tests), while progressive edge reduction yielded no significant gains on any measure. The widely held community preference for minimal-edge training (specificity argument) is not supported by the only RCT directly comparing these methods. \textit{Endpoint: dynamometric surrogate (pinch grip strength) — no on-wall performance or climbing grade outcome measured.}
\end{evidencebox}

\subsubsection{Low-Intensity Daily Hangs (``Abrahangs'')}

\begin{bioplausbox}{Observational: Low-Intensity Daily Hangs (Gilmore 2024) · \bioplausible · ESS 5}
Gilmore et al.\ \cite{GKAO24} (retrospective cohort, $n = 527$ self-selected Crimpd app users, self-reported loads) found: Abrahangs-only $+2.5\%$ strength-to-weight (\textit{not} significantly different from climbing-only, $p = 0.12$); Max Hangs-only $+3.2\%$ (significant vs.\ climbing-only, $p < 0.001$); both combined $+5.8\%$ (additive, significant vs.\ either alone). The Abrahangs-only group did not significantly outperform the climbing-only control, making a non-inferiority claim unsupported. The result is consistent with tendon mechanobiology (low-load cyclic stress promotes collagen synthesis; Baar et al.), but the design carries high risk of selection bias, confounding by total climbing volume, and effort-belief effects. Treat as hypothesis-generating only. A practical risk is that non-expert climbers may misjudge ``submaximal,'' leading to cumulative overuse.
\end{bioplausbox}

\subsubsection{Rate of Force Development}

\begin{evidencebox}{Rate of Force Development Response to Hangboarding · \evidencebacked · ESS 7}
Levernier \& Laffaye \cite{LL19} found a significant $+8\%$ RFD increase after 4 weeks of hangboarding in elite climbers. Hermans et al.\ \cite{HO22} reported a $+28\%$ within-group RFD gain in the hangboard group; however, the between-group difference was not statistically significant ($p = 0.056$, ES 0.21). The meta-analysis by Stien et al.\ \cite{SO23} pooled available trials and found a moderate-to-large effect for RFD (SDM = 0.91, 95\%~CI 0.21--1.61). Grip-depth specificity of RFD adaptation is supported by the Hermans finding that gains were larger on the trained 23~mm rung than on a jug.
\end{evidencebox}

\subsubsection{Training Volume and NMES-Assisted Loading (Dindorf et al.\ 2024)}

\begin{evidencebox}{NMES-Assisted Fingerboard Training Non-Inferior at Lower Volume (Dindorf 2024) · \evidencebacked · ESS 7}
Dindorf et al.\ \cite{DO24} randomised 18 intermediate-to-advanced climbers to twice-weekly NMES-assisted fingerboard training or three-times-weekly conventional fingerboard training for 7 weeks. Despite ${\sim}33\%$ lower training volume, the NMES group was non-inferior to the conventional group on finger flexor endurance. No significant between-group differences were observed. The finding suggests that external neuromuscular augmentation can partially compensate for reduced session frequency, with implications for athletes managing scheduling or recovery constraints. Sample size was small ($n = 18$ completers); results should be considered preliminary.
\end{evidencebox}

\subsubsection{Universal Methodological Limitations}

\begin{itemize}[noitemsep, topsep=2pt]
  \item \textbf{Sample size}: $n=11$--54 in all intervention trials (including Dindorf $n=18$); Stien et al.\ \cite{SO23} found 4 of 5 meta-analyzable trials were underpowered; PEDro scores 5--7/10 (median 6).
  \item \textbf{Cohort}: Nearly all participants are intermediate-to-advanced males; Langer et al.\ \cite{LAS24} is the only female-focused trial and found only anecdotal strength advantages.
  \item \textbf{Duration}: 4--10~weeks; tendon structural remodelling ($>\!3$~months) is not captured by any trial.
  \item \textbf{Outcome validity}: All primary outcomes are dynamometric; the Stien meta found no improvements on actual climbing-specific on-wall tests; no RCT demonstrates statistically significant grade improvement.
  \item \textbf{Confounding}: Participants continued normal climbing throughout; causal attribution is limited.
  \item \textbf{Injury data}: No RCT reports injury incidence; Sjöman et al.\ \cite{SGJ23} provides the only injury-risk data (survey-based).
\end{itemize}

\subsubsection*{Risk of Bias (Simplified Per-Study)}

\begin{longtable}{>{\raggedright\arraybackslash}p{2.8cm}
                  >{\raggedright\arraybackslash}p{1.8cm}
                  >{\raggedright\arraybackslash}p{2.0cm}
                  >{\raggedright\arraybackslash}p{2.0cm}
                  >{\raggedright\arraybackslash}p{1.6cm}
                  >{\raggedright\arraybackslash}p{1.5cm}}
\caption{Simplified risk-of-bias summary for intervention trials.}\label{tab:rob}\\
\toprule
\textbf{Study} & \textbf{Allocation conceal.} & \textbf{Participant blinding} & \textbf{Dropout / ITT} & \textbf{PEDro} & \textbf{Overall} \\
\midrule
\endfirsthead
\multicolumn{6}{c}{\textit{(continued)}}\\
\toprule
\textbf{Study} & \textbf{Allocation conceal.} & \textbf{Participant blinding} & \textbf{Dropout / ITT} & \textbf{PEDro} & \textbf{Overall} \\
\midrule
\endhead
\bottomrule
\endfoot
\bottomrule
\endlastfoot

López-Rivera 2012 & Unclear & Not feasible & Not reported & --- & High \\[2pt]
Medernach 2015 & Unclear & Not feasible & Minimal & 5 & High \\[2pt]
Levernier 2019 & Unclear & Not feasible & Not reported & 5 & High \\[2pt]
López-Rivera 2019 & Unclear & Not feasible & 32\% (38→26) & --- & High \\[2pt]
Mundry 2021 & Unclear & Not feasible & 10\% (30→27) & 6 & High \\[2pt]
Devise 2022 & Unclear & Not feasible & Not reported & 7 & Moderate \\[2pt]
Hermans 2022 & Unclear & Not feasible & Not reported & 6 & High \\[2pt]
Dindorf 2024 & Unclear & Not feasible & 18/32 enrolled (44\%) & --- & High \\[2pt]
Langer 2024 & Stated & Not feasible & 16\% (31→26) & --- & Moderate \\[2pt]
Ruiz-López 2025 & Unclear & Not feasible & 45\% (20→11) & --- & High \\[2pt]
Gilmore 2024 & N/A (obs.) & N/A & Self-selected app users & --- & High (obs.) \\[2pt]
\end{longtable}
\noindent\small PEDro scores from Stien et al.\ \cite{SO23}; `---' = not yet scored or trial postdates meta. `Not feasible' = blinding of participants impossible in exercise training trials; assessor blinding not reported in any trial. Dropout/ITT: López-Rivera 2019 analyzed 26 of 38 enrolled (32\% attrition); Ruiz-López 2025 analyzed 11 of 20 enrolled (45\% attrition) — both without ITT analysis.

\subsubsection{Meta-Analytic Evidence}

\begin{evidencebox}{Hangboard Meta-Analysis: Strength and Endurance Gains (Stien 2023) · \evidencebacked · ESS 7}
Stien et al.\ \cite{SO23} conducted a systematic review and meta-analysis (search to January 2021; 5 trials meta-analyzable, $n=110$) of hangboard and resistance training in climbers. Finger strength improved with pooled SDM~= 0.41 (95\%~CI 0.03--0.80); RFD SDM~= 0.91 (95\%~CI 0.21--1.61); forearm endurance SDM~= 1.23 (95\%~CI 0.69--1.77); overall climbing-specific test performance SDM~= 0.57 (95\%~CI 0.24--0.91). PEDro scores ranged 5--7 (median~6/10); four of five trials were underpowered ($\alpha = 0.05$, $\beta = 0.2$). \textbf{No improvement was found in actual climbing-specific on-wall tests.} \textit{Endpoint note: all pooled outcomes are dynamometric surrogate measures (force, endurance, RFD); no included trial demonstrated statistically significant improvement on a functional on-wall climbing test or grade-based outcome.} Included trials: Levernier \& Laffaye \cite{LL19}, Medernach et al.\ \cite{MKL15}, Hermans et al.\ \cite{HO22}, Mundry et al.\ \cite{MSA+21}, and Stien et al.\ (2021, campus board study). Devise et al.\ \cite{DO22}, Gilmore et al.\ \cite{GKAO24}, and 2024--2025 trials postdate the search.
\end{evidencebox}

\subsubsection{Injury and Experience-Dependent Risk (Sjöman et al.\ 2023)}

\begin{evidencebox}{Experience-Dependent Hangboard Injury Risk (Sjöman 2023) · \evidencebacked · ESS 5}
Sjöman et al.\ \cite{SGJ23} administered a web-based survey to the Finnish climbing community ($n=434$). No significant overall association between regular fingerboard training (RFT) and sport-specific pulley injury (SPIIF) was found. Stratified analysis revealed: in climbers with $\geq6$ years experience, RFT was \emph{negatively} associated with SPIIF ($\chi^2[1, n=200]=4.57$, $p=0.03$); in male climbers with $<6$ years who had reached $\geq7a$, RFT was \emph{positively} associated with SPIIF ($\chi^2[1, n=75]=4.61$, $p=0.03$). This is the only study providing quantitative injury-risk data for hangboarding, suggesting a protective effect in experienced climbers and an elevated risk in less-experienced climbers advancing rapidly.
\end{evidencebox}

\subsubsection{Female-Specific Evidence (Langer et al.\ 2024)}

The only trial enrolling exclusively female climbers \cite{LAS24} found only anecdotal evidence for fingerboard superiority over on-wall training on climbing-specific strength (BF$_{10}=1.16$). On-wall training was strongly preferred for enjoyment and intrinsic motivation. All findings in this review otherwise derive from male-advanced cohorts and should not be assumed to generalise to female or beginner populations.

%---------------------------------------------------------------------
\subsection{Practitioner Consensus and Contested Topics}
\label{sec:practitioner}

\begin{folklorebox}{Practitioner Consensus: Core Hangboard Principles · \folkloreverified · ESS 1}
Consistent across major coaching resources (\textit{Training for Climbing}, \textit{Rock Climber's Training Manual}, Lattice Training, Dave MacLeod, Crimpd, \texttt{r/climbharder} threads, 2020--2025):
\begin{enumerate}[noitemsep, topsep=2pt]
  \item \textbf{No full-bodyweight hanging for beginners below ${\sim}$1--2 years or ${\sim}$V4/5.11} (traditional majority view): technique development is more efficient and tendons are not yet conditioned for isolated peak loads. \textit{Contested: a significant and growing camp — including Crimpd, Emil Abrahamsson (creator of Abrahangs, targeting climbers at ${\sim}5.11$), Lattice Training's feet-on-floor progressions, and Tyler Nelson's submaximal tissue-tolerance protocols — explicitly endorses early-stage low-intensity loading (feet on floor, $\leq40\%$ MFS) to build baseline tendon capacity before progressing to full bodyweight. The shift gained momentum on \texttt{r/climbharder} post-2022 and is reflected in current app-based training plans. The key distinction is ``no max-intensity loading'' vs.\ ``no finger loading whatsoever'': the former is near-universal consensus, the latter is increasingly contested.}
  \item \textbf{Minimum 15-min warm-up}: light cardio, shoulder/wrist mobility, and progressive large-hold hangs before any session.
  \item \textbf{Session freshness required}: hangboard before climbing or on a separate day; never after exhausting bouldering.
  \item \textbf{2--3 sessions/week with $\geq48$~h between hard sessions}. \textit{Note: this directly conflicts with Abrahamsson's twice-daily low-intensity protocol (Abrahangs), discussed under Contested Topics.}
  \item \textbf{Half-crimp or open-hand as default}: full crimp introduced only gradually and specifically.
  \item \textbf{Progress by load before reducing edge depth}: overload an 18--20~mm edge before moving to smaller holds.
  \item \textbf{Deload every 4--8~weeks}: reduce volume or intensity ${\sim}40$--50\% for one week to allow structural adaptation.
\end{enumerate}
\end{folklorebox}

\begin{bioplausbox}{Early Max-Intensity Loading Risk in Novice Climbers · \bioplausible · ESS 2}
Point 1 is clinically plausible---finger tendons adapt more slowly than muscles, and isolated peak loading before adequate conditioning likely elevates A2 rupture risk---and receives partial empirical support from Sjöman et al.\ \cite{SGJ23}, who found RFT positively associated with injury in less-experienced climbers reaching 7a$+$ ($p=0.03$). No controlled study has measured injury rates in beginner hangboard users vs.\ beginner-only climbers; the consensus threshold of 1--2~years / V4 is expert opinion, not an empirically derived value.
\end{bioplausbox}

\subsubsection*{Contested Topics}

\textbf{20~mm as the standard edge.} The community treats a 20~mm half-crimp as the default test and training edge on the grounds that it isolates the finger flexors without excessive skin-friction contribution yet still accommodates progressive loading. No controlled study directly compares edge depths as training stimuli; the recommendation is expert consensus, not an empirically derived threshold.

\textbf{Minimal-edge vs.\ added-weight training.} Community preference for minimal-edge work (specificity argument: hard climbing occurs on small holds) directly conflicts with Mundry et al.\ \cite{MSA+21}, the only RCT comparing the two methods, which found added weight on a standard edge significantly improved grip strength while edge reduction did not.

\textbf{Daily low-intensity hangs.} Abrahamsson's twice-daily submaximal protocol conflicts with the 48~h recovery rule but is mechanobiologically plausible and observationally supported \cite{GKAO24}. No RCT has tested this modality directly.

\textbf{No-hang devices.} No-hang tools (Tension Block, Grippul) offer precise weight-based loading with full finger-flexor isolation and no shoulder involvement. No controlled study compares them to traditional hangboarding; the clinical significance of absent shoulder engagement is unresolved.

\textbf{Full-crimp training.} Traditional guidance avoids full-crimp hangboarding given the high pulley loads \cite{Sch01}. Contemporary practitioners increasingly argue for progressive full-crimp exposure when it is used extensively on the wall, reasoning that tissue-specific inoculation is necessary. No RCT has tested this.

\input{sections/discussion}

%---------------------------------------------------------------------
\section{Weekly Training Integration}
\label{sec:weekly}

%---------------------------------------------------------------------
\subsection{Summary of Cited Peer-Reviewed Studies}
\label{subsec:weekly_table}

\noindent\textbf{\evidencebacked} studies providing direct quantitative data for this section are listed below. Studies not independently confirmed by numeric effect-size verification are excluded from the table and flagged where used in text.

{\footnotesize
\setlength{\tabcolsep}{4pt}
\renewcommand{\arraystretch}{1.2}
\begin{longtable}{>{\raggedright\arraybackslash}p{2.2cm}
                  >{\raggedright\arraybackslash}p{2.0cm}
                  >{\raggedright\arraybackslash}p{2.6cm}
                  >{\raggedright\arraybackslash}p{2.5cm}
                  >{\raggedright\arraybackslash}p{2.8cm}
                  >{\raggedright\arraybackslash}p{1.4cm}
                  >{\raggedright\arraybackslash}p{1.2cm}}
\caption{Peer-reviewed studies informing weekly training integration}
\label{tab:weekly_studies}\\
\toprule
\textbf{Authors} & \textbf{Design / $n$} & \textbf{Protocol} & \textbf{Primary outcome} & \textbf{Effect size} & \textbf{Endpoint} & \textbf{Quality} \\
\midrule
\endfirsthead
\multicolumn{7}{l}{\small\itshape Table \ref{tab:weekly_studies} continued\ldots} \\
\toprule
\textbf{Authors} & \textbf{Design / $n$} & \textbf{Protocol} & \textbf{Primary outcome} & \textbf{Effect size} & \textbf{Endpoint} & \textbf{Quality} \\
\midrule
\endhead
\bottomrule
\endlastfoot

Stien et al.\ \cite{SO23}
& SR + meta-analysis; $n = 11$ studies (5 pooled RCTs)
& Climbing resistance training 2–3$\times$/wk, 4–10 wk vs.\ climbing-only control
& Finger strength, RFD, forearm endurance (dynamometric)
& Overall SDM = 0.57 (95\% CI 0.24–0.91); strength 0.41 (0.03–0.80); RFD 0.90 (0.21–1.61); endurance 1.23 (0.69–1.77)
& Surrogate; no on-wall gain
& GRADE: low–moderate \\

López-Rivera \& González-Badillo \cite{LRGB19}
& RCT, 3 groups; $n = 26$ advanced ($\sim$7c+/8a)
& 8 wk, 2$\times$/wk: MaxHangs ($\geq$80\% MFS) vs.\ IntHangs vs.\ Combination; all climbed 6 d/wk
& Grip endurance: time-to-failure dead hang on 11 mm edge
& IntHangs $+$45\% (ES 1.0, $p=0.002$); MaxHangs $+$34\% (ES 0.6, $p=0.010$); Combo $+$7\% (ns)
& Surrogate (dynamometric)
& GRADE: low \\

Devise et al.\ \cite{DO22}
& RCT, 4 groups; $n = 54$ experienced (41M/13F)
& 4 wk, 2$\times$/wk: F100 / F80 / F60 / control on 12 mm hold
& MFS (N), stamina, endurance
& F100 MFS $+$14.3\% ($p<0.001$); F80 MFS $+$12.4\%; F80 stamina 29.5$\to$60.4\%; F80 endurance 63.2$\to$79.2\%
& Surrogate (dynamometric)
& GRADE: low \\

Stien et al.\ \cite{SPV+21}
& RCT; $n = 16$ adv–elite (groups $n=5,6,5$)
& 5 wk, campus board: 4$\times$/wk vs.\ 2$\times$/wk vs.\ climbing-only; $\geq$48 h between sessions
& Bouldering performance, max force, campus moves, RFD
& 4$\times$/wk: force/RFD $\uparrow$ vs.\ control; 2$\times$/wk: bouldering/campus $\uparrow$ vs.\ control; no between-freq.\ difference
& Mixed (bouldering: functional; force/RFD: surrogate)
& GRADE: very low \\

Wilson et al.\ \cite{WMR+12}
& Meta-analysis; 21 RCTs/CCTs; $n = 422$ ES
& Concurrent strength + endurance vs.\ strength-only; running/cycling; varied freq.\ and duration
& Hypertrophy, maximal strength, power
& Strength ES: 1.76 (solo) vs.\ 1.44 (concurrent, $-$18\%); hypertrophy $-$31\%; power $-$40\%; endurance duration $r=-0.29$ to $-0.75$
& Surrogate; non-climbing; inferential transfer
& GRADE: moderate \\

\end{longtable}
}% end footnotesize

%---------------------------------------------------------------------
\subsection{Fatigue Interference Between Modalities}
\label{subsec:fatigue_interference}

\begin{evidencebox}{Supplemental Finger Training Outperforms Climbing Alone (Stien 2023 meta-analysis) · \evidencebacked · ESS 7}
Supplemental finger-strength training (hangboard or campus board) superimposed on regular climbing training produces significantly greater gains in climbing-specific tests than climbing alone: overall meta-analytic SDM = 0.57 (95\% CI 0.24–0.91), finger strength SDM = 0.41 (0.03–0.80), forearm endurance SDM = 1.23 (0.69–1.77) (Stien et al.\ 2023, \textit{Biol Sport} 40(1):179–191; \cite{SO23}). \textit{Endpoint: all pooled outcomes are dynamometric surrogates (force, endurance, RFD); no included trial showed significant improvement on an on-wall climbing test.} These studies implicitly impose $\geq$48 h between sessions by scheduling 2–3 sessions/wk but do not experimentally vary session order or inter-session gap. No published RCT has directly randomized the temporal proximity between hangboard training and limit bouldering to measure interference on adaptation outcomes in climbers.

The concurrent-training interference effect (Wilson et al.\ 2012) is well-characterised for lower-body strength plus cardiovascular endurance: adding endurance training at high frequency ($r = -0.35$) or duration ($r = -0.75$) significantly attenuates strength, hypertrophy, and power gains versus strength-only training, with running producing larger deficits than cycling. The molecular basis is AMPK-mediated inhibition of mTORC1: AMPK requires $\geq$3 h post-endurance to return to baseline, while mTORC1 remains elevated $\approx$18 h post-resistance loading.
\end{evidencebox}

\begin{bioplausbox}{Concurrent Training Interference: Peripheral and Molecular Mechanisms · \bioplausible · ESS 2}
Limit bouldering ($\geq$90–100\% effort) and hangboard max-hang sessions share the same primary effector tissue: the finger flexor musculotendinous unit (FDP/FDS) and annular pulley system (A2, A4). Both impose near-maximal isometric loading on a 6–14 mm crimp, activating the same motor-unit pool at high discharge rates. The dominant interference mechanism is not the classical AMPK/mTOR competition (which is driven by systemic aerobic endurance), but rather: (1) acute peripheral fatigue — H\textsuperscript{+} accumulation and impaired sarcoplasmic Ca\textsuperscript{2+} release reduce peak force capacity; (2) phosphocreatine depletion — PCr returns to $\approx$50\% of resting levels within 2–4 min but full resynthesis takes $\approx$8–12 min post-maximal isometric effort; (3) central inhibition — reduced corticospinal drive via Golgi tendon organ (GTO) afference after maximal isometric loading. A second max-effort session initiated before full neuromuscular recovery delivers a degraded force stimulus, reducing the adaptive signal for maximal strength.

The AMPK/mTOR model applies most cleanly to combining route climbing endurance (aerobic forearm demand) with hangboard or limit bouldering (maximal isometric). Climbing endurance is a local aerobic-to-anaerobic continuum without the eccentric locomotion that drives the strongest interference in Wilson et al.\ (2012), so the magnitude of molecular interference in climbers is likely $>0$ but smaller than the $-$18\% strength deficit observed in running-concurrent protocols.

\textit{Evidence gap: no RCT has measured post-hangboard neuromuscular recovery time in the finger flexors at a temporal resolution sufficient to guide programming, and no study has measured AMPK or mTOR pathway activity in finger flexors following climbing-specific endurance versus limit bouldering.}
\end{bioplausbox}

\begin{folklorebox}{48\,h Minimum Between High-Intensity Finger Sessions · \folkloreverified · ESS 1}
\textbf{Source: Lattice Training blog (Randall \& Torr); Anderson \& Anderson, \textit{The Rock Climber's Training Manual}, 2nd ed., 2014; Hörst, \textit{Training for Climbing}, 3rd ed., 2016; r/climbharder wiki.}

Practitioner consensus across all major coaching frameworks prescribes $\geq$48 h between any two sessions imposing high-intensity finger loading (max-hangs, campus board, limit bouldering, Moonboard benchmarks). Some coaches (Neil Gresham, UKClimbing; Power Company Podcast) recommend 72 h between max-intensity finger sessions for advanced climbers. Both the 48 h and 72 h thresholds are study-design conventions embedded in climbing RCTs as methodological safeguards (Stien 2021 specified $\geq$48 h between campus sessions), not empirically optimised prescriptions.
\end{folklorebox}

\subsubsection*{Moonboard Sessions}

\begin{bioplausbox}{Moonboard as High-Intensity Finger Session: Scheduling Equivalence · \bioplausible · ESS 2}
The Moonboard is a 40$^\circ$-overhung standardised board using benchmark problems with $\leq$1-phalange-depth holds. A full session (60–90 min, $\geq$8–10 max-effort attempts per problem) involves near-maximal finger loading, high CNS demand, and significant metabolic depletion — a fatigue profile overlapping limit bouldering on a steep wall, not top-rope endurance. Moonboard sessions should be categorised as high-intensity finger sessions equivalent to limit bouldering for scheduling purposes; the $\geq$48 h separation constraint applies identically. Placing a Moonboard session adjacent to a hangboard max-strength day (e.g., Moonboard Monday, Hangboard Tuesday) risks the same attenuation of MFS stimulus quality as doing limit bouldering the day before max-hangs.

\textit{Evidence gap: no study has characterised post-Moonboard recovery of MFS, RFD, or tendon stiffness, nor compared its fatigue profile to an equivalent-duration limit-bouldering session.}\singlesource
\end{bioplausbox}

\subsubsection*{Top-Rope and Lead Endurance Sessions}

\begin{bioplausbox}{Route Endurance Fatigue Profile: Lower Interference Than Limit Bouldering · \bioplausible · ESS 2}
Top-rope and lead climbing at 60–75\% maximum effort impose forearm endurance demand (aerobic-glycolytic) with substantially lower peak finger force per hold than limit bouldering (estimated 50–70\% MFS vs.\ $\geq$90\% MFS). Active recovery research in climbers demonstrates that low-to-moderate climbing as an inter-session recovery modality accumulates forearm fatigue relative to non-climbing locomotion (performance impulse 29.5\% lower after climbing vs.\ 18.4\% lower after walking, $p = 0.009$, $d = 0.54$).\singlesource The fatigue profile of route endurance (sustained aerobic demand at submaximal grip force) is qualitatively distinct from max-hang fatigue and is unlikely to degrade max-strength stimulus quality as severely as placing a Moonboard session adjacent to a hangboard day. Optimal placement: route endurance sessions should bookend the high-intensity block, either as an early-week warm-up modality or as the final training day before a rest day. They should \emph{not} immediately precede a hangboard max-strength session.
\end{bioplausbox}

%---------------------------------------------------------------------
\subsection{Session Sequencing Rules}
\label{subsec:sequencing}

\begin{bioplausbox}{Hangboard Before Bouldering for Max-Strength Priority · \bioplausible · ESS 2}
\textbf{Hangboard before or after limit bouldering?}

Devise et al.\ (2022) demonstrate intensity-specific adaptation: F100 (100\% MFS) significantly improves MFS but not stamina or endurance; F80 (80\% MFS) improves MFS, stamina, and endurance; F60 (60\% MFS) improves stamina and endurance but not MFS. This specificity has direct sequencing implications: max-hang adaptation requires delivery of the stimulus at close to peak force output. Any prior modality that depletes finger-flexor force capacity shifts the effective stimulus downward — an intended F100 session becomes an F80 or F60 session after a 2–3 h limit bouldering session (estimated MVC decrement 10–25\%, inferred from general post-climbing strength decrements; no direct hangboard-after-bouldering measurement exists in the peer-reviewed literature).

\textbf{Principle:} When the goal is maximal finger strength or RFD, hangboard max-hangs must precede — not follow — limit bouldering, and should be executed at session start after warm-up. When the goal is endurance (repeaters at F60–F80 intensity), this constraint is weaker because the submaximal stimulus is less sensitive to pre-fatigue.

\textit{Evidence gap: no RCT has randomised hangboard-before vs.\ hangboard-after bouldering to directly test the sequence effect in climbers.}
\end{bioplausbox}

\begin{bioplausbox}{Same-Day Hangboard + Limit Bouldering: Feasibility Conditions · \bioplausible · ESS 2}
\textbf{Same-day hangboard + limit bouldering.}

Same-session combination is feasible if: (a) max-strength hangboard is performed first, immediately after warm-up; (b) within-session passive gap of approximately 10--25\,min (broad range; no climbing data; general PCr resynthesis literature suggests $\approx$8--12\,min for full restoration post-maximal isometric effort) provides partial PCr recovery without additional forearm loading; (c) total hangboard volume reduced approximately 30--60\% versus a dedicated session (a pragmatic range reflecting the absence of dose-response climbing data; individual tolerance varies); (d) total high-intensity time-under-tension ideally limited to roughly 45--90\,s across both modalities to moderate AMPK activation (these values are ESS~2 extrapolations from concurrent-training literature, not empirically derived thresholds for climbing). Combining both modalities in one session does not substitute for two separated sessions, as cumulative volume limits the training signal of each.\singlesource

Sequential strength-before-endurance within a concurrent session preserves the quality of the strength stimulus in general concurrent-training literature; the same principle applies by bioplausible transfer to upper-limb isometric loading.
\end{bioplausbox}

\begin{folklorebox}{Hangboard-First Session Ordering · \folkloreverified · ESS 1}
\textbf{Source: Eva López periodization blog; Lattice Training blog; Anderson \& Anderson 2014.}

Eva López's periodization guidelines explicitly recommend hangboard work at the \emph{start} of a session when freshness is required for max-effort protocol quality. The Anderson brothers and Lattice Training both advocate same-day hangboard-first ordering when schedule constraints prevent separate days. These are practitioner heuristics without RCT validation.
\end{folklorebox}

%---------------------------------------------------------------------
\subsection{Example Weekly Structures}
\label{subsec:templates}

\noindent\textbf{Scope:} All templates target intermediate-to-advanced climbers (V6–V9 / 5.12–5.13c redpoint), 4–5 training days/wk. Intensity descriptors: \textit{max} = $\geq$90–100\% MFS; \textit{high submaximal} = 75–89\% MFS; \textit{submaximal} = 60–74\% MFS; \textit{technique/recovery} = $\leq$60\%, skill-focused. No complete 7-day weekly structure has been tested as a unit in a climbing-specific RCT; component protocols carry their individual evidence ratings.

\subsubsection*{Template A — Strength-Priority Block (4 days/wk, hangboard-led)}

\noindent\textit{Primary goal: maximal finger strength and RFD. Most directly evidence-proximate template; hangboard protocol matches the 2$\times$/wk frequency used in López-Rivera \& González-Badillo (2019) and consistent with Devise et al.\ (2022) F100/F80 parameters.}

\begin{itemize}[leftmargin=1.6cm,labelwidth=1.2cm,align=left]
  \item[\textbf{Mon:}] \textbf{Hangboard max-hangs.} $\geq$80\% MFS, 10 s hangs, 3–5 sets $\times$ 3–5 reps, 3–5 min inter-set rest, 18–20 mm edge; total $\approx$20–30 min. \textbf{Evidence-backed:} López-Rivera \& González-Badillo (2019) — MaxHangs protocol structure and intensity; Devise et al.\ (2022) — F100/F80 intensity parameters.
  \item[\textbf{Tue:}] \textbf{Rest or non-climbing antagonist work} (wrist extensors, shoulder external rotation), $\leq$30 min. \textbf{Bioplausible:} passive recovery to allow neuromuscular restoration; 24 h post-hangboard likely insufficient for full MVC recovery.
  \item[\textbf{Wed:}] \textbf{Limit bouldering / Moonboard projecting.} 90–100\% effort; 2–4 attempts per project problem; 5–10 min rest between attempts; climbing time $\approx$60–90 min; total session $\approx$2–2.5 h including warm-up. \textbf{Folklore:} Anderson \& Anderson 2014, r/climbharder wiki — structure of projecting sessions; no RCT specifies volume/rest parameters for limit bouldering as a training modality.
  \item[\textbf{Thu:}] \textbf{Rest.} $\geq$48 h separation from Wednesday's high-intensity finger session before next max-load day. \textbf{Bioplausible.}
  \item[\textbf{Fri:}] \textbf{Hangboard max-hangs} (second session; same as Monday, or Moonboard benchmark problems at 75–85\% effort, $\approx$90 min if goal shifts to power-endurance). \textbf{Evidence-backed} for hangboard frequency (2$\times$/wk) per Stien et al.\ (2021) and López-Rivera \& González-Badillo (2019). Moonboard substitution: \textbf{Folklore} (Lattice Training periodization framework).
  \item[\textbf{Sat:}] \textbf{Top-rope/lead endurance.} Routes at 60–75\% max effort; $\approx$2 h; $\geq$10 routes or equivalent mileage. \textbf{Bioplausible:} aerobic forearm demand stimulates capillarisation and lactate-clearance adaptations; low enough intensity not to significantly attenuate next-week hangboard quality.
  \item[\textbf{Sun:}] \textbf{Rest.}
\end{itemize}

\noindent\textit{Classification: Hangboard days are the most directly evidence-backed component. The overall 7-day arrangement is folklore-derived; the sequencing logic (hard finger $\to$ rest $\to$ limit bouldering $\to$ rest $\to$ hangboard $\to$ endurance $\to$ rest) is bioplausible from concurrent-training principles.}

\subsubsection*{Template B — Bouldering-Led Block (4 days/wk, board-climbing priority)}

\noindent\textit{Primary goal: project-specific bouldering performance and on-wall power. Hangboard is maintained at reduced volume as a supplemental stimulus.}

\begin{itemize}[leftmargin=1.6cm,labelwidth=1.2cm,align=left]
  \item[\textbf{Mon:}] \textbf{Limit bouldering.} 90–100\% effort; project-style attempts; 2–2.5 h total (30 min warm-up, 60–90 min projecting). \textbf{Folklore:} Anderson \& Anderson 2014 ``Limit Bouldering'' chapter; r/climbharder wiki.
  \item[\textbf{Tue:}] \textbf{Hangboard max-hangs (reduced volume).} $\geq$80\% MFS; 10 s, 3 $\times$ 4–5 reps; performed 24–48 h post-bouldering — a scheduling compromise; volume reduced 20–30\% vs.\ a fully recovered session if within 48 h. \textbf{Bioplausible:} sub-optimal recovery window; \textbf{Folklore:} Lattice Training blog explicitly flags the 24 h post-bouldering hangboard trade-off.
  \item[\textbf{Wed:}] \textbf{Rest or antagonist/mobility.} \textbf{Bioplausible.}
  \item[\textbf{Thu:}] \textbf{Moonboard benchmark problems.} 70–85\% effort; emphasis on movement quality over max-effort; $\approx$90 min. \textbf{Folklore:} Moonboard app programming, r/climbharder.
  \item[\textbf{Fri:}] \textbf{Rest.} $\geq$48 h before next high-intensity day. \textbf{Bioplausible.}
  \item[\textbf{Sat:}] \textbf{Top-rope/lead endurance or ARC.} 60–75\% max effort, $\approx$2 h. \textbf{Folklore:} Anderson \& Anderson 2014 — ARC training definition and prescription.
  \item[\textbf{Sun:}] \textbf{Rest.}
\end{itemize}

\noindent\textit{Classification: Folklore-derived overall structure; individual hangboard parameters are evidence-backed. The Tuesday hangboard after Monday bouldering is a pragmatic scheduling compromise, classified as bioplausible rather than optimal.}

\subsubsection*{Template C — Evidence-Constrained Balanced Block (4 days/wk)}

\noindent\textit{Primary goal: balanced finger-strength and endurance development. This template most closely replicates the 2$\times$/wk hangboard-plus-normal-climbing design used across Hermans et al.\ (2022)\footnote{Hermans et al.\ (2022) \textit{Front Sports Act Living} 4:888158. \cite{HO22}. RCT; $n = 35$; 10 wk, 2 sessions/wk (48 min/session), Beastmaker 1000 app-guided protocol; $F_\mathrm{peak}$ $\uparrow$ ($p < 0.05$, large within-group ES) vs.\ control. GRADE: low. \toverify: cited by fewer than 3 independent sources; requires verification. See Appendix~\ref{app:toverify}, entry~2.} and consistent with López-Rivera \& González-Badillo (2019).}

\begin{itemize}[leftmargin=1.6cm,labelwidth=1.2cm,align=left]
  \item[\textbf{Mon:}] \textbf{Hangboard protocol} (2$\times$/wk schedule; Beastmaker-style or equivalent force-quantified protocol; 48 min including warm-up; progressive overload weekly). \textbf{Evidence-backed:} 2$\times$/wk frequency consistent with López-Rivera \& González-Badillo (2019) and Devise et al.\ (2022) intervention structures.
  \item[\textbf{Tue:}] \textbf{Rest.} $\geq$48 h minimum between hangboard sessions. \textbf{Evidence-backed:} study constraint in climbing-specific RCTs (Stien 2021; Devise 2022).
  \item[\textbf{Wed:}] \textbf{Limit bouldering or Moonboard.} 90–100\% effort; 60–90 min climbing; equivalent to ``normal climbing'' continued by participants in the intervention arms of the cited RCTs. \textbf{Folklore} for specific volume/intensity; consistent with the design intent of the cited trials.
  \item[\textbf{Thu:}] \textbf{Rest or antagonist work.} \textbf{Bioplausible.}
  \item[\textbf{Fri:}] \textbf{Hangboard protocol} (second session of week; identical structure to Monday). \textbf{Evidence-backed:} 2$\times$/wk schedule per intervention design.
  \item[\textbf{Sat:}] \textbf{Route climbing/lead endurance.} 60–75\% max effort; $\approx$2 h. \textbf{Bioplausible:} aerobic demand modality separate from finger-strength focus days.
  \item[\textbf{Sun:}] \textbf{Rest.}
\end{itemize}

\noindent\textit{Classification: The hangboard days are the most directly evidence-backed element of any template. The specific 7-day composition is folklore-derived, but internally consistent with the evidence-based principle of $\geq$48 h between high-intensity finger sessions.}

\subsubsection*{General Constraint Across All Templates}

\begin{folklorebox}{Weekly Hard Finger Exposure Cap: 2--3 Sessions · \folkloreverified · ESS 1}
\textbf{Source: Lattice Training; Anderson \& Anderson 2014; r/climbharder wiki.}

Keep total ``hard finger exposures'' (limit bouldering, Moonboard, max-hangs at $\geq$80\% MFS) at \textbf{2–3 per week} for most V6–V9 climbers unless individual recovery capacity has been demonstrated to support higher frequency. This cap is a widespread coaching heuristic with no RCT dose-response validation in the target population.
\end{folklorebox}

%---------------------------------------------------------------------
\subsection{Deload Triggers and Protocol}
\label{subsec:deload}

\begin{bioplausbox}{Overreaching Taxonomy and Deload Indicators for Climbers · \bioplausible · ESS 2}
\textbf{Overreaching taxonomy (from non-climbing sports literature; bioplausible transfer).}

Functional overreaching (FOR) = planned short-term performance decrement with supercompensation on recovery, full restoration $\leq$2 wk. Non-functional overreaching (NFOR) = persistent decrements requiring $\leq$4 wk to resolve. Overtraining syndrome (OTS) = months-to-years recovery. FOR is tolerated training stress; NFOR is the injury-adjacent state. Resting HR elevation of 3–5 bpm above baseline with HRV recovering within 24–72 h after a rest day indicates FOR; persistent HR elevation $>$5–10 bpm above baseline on rest days with declining HRV over $\geq$5 consecutive training days indicates NFOR.

Deload indicators for climbing application (ESS~2 — all values extrapolated from general resistance-training literature; none validated in climbing-specific studies): (1) max-hang force output declining roughly $>$5--15\% over 2--3 consecutive sessions vs.\ 4-week rolling average (a plausible signal range; the exact threshold of ``10\%'' carries false precision and should be treated as an approximate sentinel); (2) subjective RPE for a fixed hangboard protocol increasing $\geq$1--2 points (0--10 scale) without load increase (this range reflects inter-individual RPE variability; a single 1.5-point increase is insufficient signal in isolation); (3) finger/pulley pain on VAS $\geq$3/10 at rest or $\geq$5/10 during loading — \textbf{mandatory hard stop} regardless of other indicators (may signal A2/A4 pulley irritation requiring medical evaluation; this threshold is the one most directly supported by injury-risk rationale, even though no RCT has validated the specific VAS cut-point); (4) sustained HRV (RMSSD) downward trend over $\geq$4--7 consecutive training days combined with any two of markers (1), (2), or (5) (the ``5 days'' figure is a convention from endurance coaching, not a validated climbing threshold); (5) warm-up grades feeling $\geq$1 full grade harder than baseline for $\geq$5--10 days.

Connective tissue (pulley, tendon) adapts more slowly than muscle ($\approx$12 wk for collagen turnover); neuromuscular indicators may lag soft-tissue overload, making subjective fatigue an imperfect sentinel for pulley injury risk.

\textit{Evidence gap: no climbing-specific study has validated any HRV threshold, RMSSD trajectory, or RPE criterion for deload timing in climbers; the thresholds above are extrapolated from general resistance-training literature.}
\end{bioplausbox}

\begin{bioplausbox}{Deload Protocol: 40--60\,\% Volume Reduction for 7 Days · \bioplausible · ESS 2}
\textbf{Deload protocol.} Reduce total volume 40–60\% across all modalities simultaneously for 7 days. Maintain $\geq$1 high-intensity exposure (e.g., 2–3 sets of max-hangs at $\geq$80\% MFS) to preserve neural qualities. Eliminate repeated-failure pump sessions and any session performed to forearm failure. Simultaneous reduction of all modalities restores AMPK/mTOR balance faster than isolated forearm rest when route climbing has been part of the loading block.
\end{bioplausbox}

\begin{folklorebox}{Planned Deload Every 3--6 Weeks · \folkloreverified · ESS 1}
\textbf{Source: Lattice Training (Progression Pyramid framework); Anderson \& Anderson 2014; Hörst 2016; r/climbharder wiki.}

Practitioner consensus recommends a planned deload (50–60\% volume reduction, maintained intensity $\geq$70\% MFS) every 3–6 weeks: Lattice schedules a recovery week after every 3-week loading block; the Anderson brothers specify a rest week every 4th week. Additional folklore triggers: Moonboard benchmark send rate halving for 2 consecutive weeks at the same effort level; inability to complete $>$80\% of planned hard moves for 2 consecutive sessions. Neither schedule nor trigger threshold has been tested in a climbing-specific RCT.
\end{folklorebox}

%---------------------------------------------------------------------
\subsection{What the Literature Cannot Yet Answer}
\label{subsec:open_questions}

The following gaps constitute genuine absences in the peer-reviewed climbing literature as of April 2026.

\begin{enumerate}

  \item \textbf{Optimal inter-session recovery interval.} No RCT has compared 24 h vs.\ 48 h vs.\ 72 h spacing between two high-intensity finger sessions ($\geq$80\% MFS hangboard or limit bouldering) in trained climbers, measuring MFS, RFD, and tendon morphology (ultrasound cross-section) as endpoints. The 48 h heuristic embedded in study designs is a methodological convention, not an empirically optimised prescription.

  \item \textbf{Session sequencing effect on long-term adaptation.} No RCT has randomised hangboard-before vs.\ hangboard-after bouldering (same day or adjacent days) and measured 8–12 week finger-strength or climbing-performance outcomes with matched total weekly volume. All current sequencing rules are extrapolated from concurrent-training principles and practitioner experience.

  \item \textbf{Moonboard-specific fatigue profile and interference.} No study has characterised post-Moonboard recovery of MFS or tendon stiffness, quantified the interference between Moonboard sessions and hangboard max-hangs within the same microcycle, or compared the Moonboard's fatigue profile to equivalent-duration limit bouldering on a less-steep wall.

  \item \textbf{Concurrent-training interference magnitude in climbing.} No RCT has measured AMPK or mTOR pathway activity in finger flexors following route-climbing endurance versus limit bouldering, nor directly compared adaptation outcomes in climbers who separate these modalities by 6 h, 24 h, or 48 h. The application of Wilson et al.\ (2012) to isometric forearm loading remains inferential.

  \item \textbf{Maintenance frequency during redpointing phases.} Stien et al.\ (2021) compared 2$\times$/wk vs.\ 4$\times$/wk campus board for acquisition; no study has examined 1$\times$/wk vs.\ 2$\times$/wk maintenance of finger-strength adaptations during a concurrent limit-bouldering and route-climbing performance phase.

  \item \textbf{Validated readiness and deload-timing metrics.} No climbing-specific study has validated any HRV threshold, performance-decrement criterion, or multivariate readiness algorithm that reliably predicts functional overreaching or pulley-injury risk in V6–V9 climbers undergoing concurrent modality programmes. Individual pulley load tolerance likely varies substantially by training age, sex, and tissue stiffness — moderating variables that have not been systematically quantified.

\end{enumerate}

%---------------------------------------------------------------------
\section{Supplementation}
\label{sec:supplements}

%---------------------------------------------------------------------
\subsection*{Overview}

No supplement covered in this section has been tested in a climbing-specific RCT with a functional primary outcome (redpoint grade, route completion, pulley injury incidence, or return-to-climbing time). All evidence tiers for climbing applications are therefore extrapolated from adjacent populations or from surrogate-biomarker protocols. Claims confirmed across multiple independent literature queries are reported without qualification; single-source claims carry an explicit footnote and a \toverify flag where applicable.

%---------------------------------------------------------------------
\subsection{Study Summary Table}
\label{subsec:supptable}

\begin{longtable}{>{\raggedright\arraybackslash}p{2.8cm}
                  >{\raggedright\arraybackslash}p{2.4cm}
                  >{\raggedright\arraybackslash}p{2.6cm}
                  >{\raggedright\arraybackslash}p{2.6cm}
                  >{\raggedright\arraybackslash}p{2.0cm}
                  >{\raggedright\arraybackslash}p{1.6cm}
                  >{\raggedright\arraybackslash}p{1.0cm}
                  >{\raggedright\arraybackslash}p{1.2cm}}
\caption{Peer-reviewed intervention studies cited in this section.
         P1NP = amino-terminal propeptide of type I procollagen (collagen synthesis surrogate).
         VISA-A = Victorian Institute of Sport Assessment -- Achilles tendinopathy score.
         OSS = Oxford Shoulder Score; SPADI = Shoulder Pain and Disability Index.
         g = Hedges $g$; $d$ = Cohen's $d$.
         Endpoint: S = surrogate; F = functional.
         ES not reported = authors did not publish standardised effect size or 95\,\% CI.}\label{tab:suppstudy}\\
\toprule
\textbf{Study} & \textbf{Design / n} & \textbf{Supplement / dose} & \textbf{Primary outcome} & \textbf{Effect size} & \textbf{Endpoint} & \textbf{Qual.} & \textbf{Industry} \\
\midrule
\endfirsthead
\toprule
\textbf{Study} & \textbf{Design / n} & \textbf{Supplement / dose} & \textbf{Primary outcome} & \textbf{Effect size} & \textbf{Endpoint} & \textbf{Qual.} & \textbf{Industry} \\
\midrule
\endhead
\midrule
\multicolumn{8}{r}{\small\itshape Continued on next page}\\
\endfoot
\bottomrule
\endlastfoot

Shaw et al.\ (2017).
\textit{Am J Clin Nutr} 105(1):136--143. \cite{SLBR+17}
&
Randomised double-blind crossover; n=8 healthy males, age $21{\pm}0.5$\,yr; no tendon pathology. Washout 7\,days.
&
5\,g or 15\,g gelatin $+$ 48\,mg vitamin~C vs.\ placebo; single dose 1\,h before 6\,min rope-skipping.
&
Serum P1NP (AUC, 4\,h); mechanical properties of in-vitro engineered ligament constructs incubated with participant serum.
&
15\,g: ${\approx}2{\times}$ P1NP AUC vs.\ placebo ($p{<}0.05$); in-vitro ligament stiffness and collagen content $\uparrow$ significantly. Formal Cohen's $d$ / 95\,\%\,CI not reported in primary publication.%
\footnote{Single-source: P1NP at 4\,h reported as $+153\,\%$ vs.\ baseline for 15\,g and $+53.9\,\%$ for placebo; exact source tables not independently verified.}
&
\textbf{S} (biomarker P1NP; in-vitro ligament proxy; no tendon structural or clinical outcome)
&
Low
&
Baar lab; gelatin donated by industry; co-author holds related patents. Flag.\\

\addlinespace

Lis \& Baar (2019).
\textit{Int J Sport Nutr Exerc Metab} 29(5):526--531. \cite{LB19}
&
Randomised double-blind crossover; n=10 recreationally active males.
&
15\,g vitamin-C-enriched gelatin, 15\,g vitamin-C-enriched hydrolysed collagen, gummy (gelatin+HC), or placebo; 1\,h before 6\,min jump rope.
&
Serum P1NP at 4\,h; serum amino-acid profile.
&
Gelatin and HC both: ${\approx}20\,\%$ rise in P1NP from baseline (trend), not statistically significant vs.\ placebo due to high inter-individual variability. ES and 95\,\%\,CI not reported.
&
\textbf{S} (biomarker P1NP; amino-acid kinetics; no structural or clinical endpoint)
&
Low
&
Same lab as Shaw (2017); no explicit conflict statement.\\

\addlinespace

Praet et al.\ (2019).
\textit{Nutrients} 11(1):76. \cite{PPW+19}
&
Double-blind placebo-controlled crossover pilot RCT; n=20 chronic mid-portion Achilles tendinopathy patients; 6-month protocol.
&
5\,g/day specific collagen peptides (TENDOFORTE\textsuperscript{\textregistered}) $+$ twice-daily calf strengthening vs.\ placebo $+$ strengthening.
&
VISA-A score (pain/function) at 3\,months; tendon microvascularity (Doppler).
&
VISA-A: collagen-first $+12.6$ (95\,\%\,CI 9.7--15.5) vs.\ placebo-first $+5.3$ (2.3--8.3). MCID for VISA-A ${\approx}6.5$. Microvascularity decreased similarly in both groups (no between-group difference).
&
\textbf{F} (VISA-A: validated pain/function scale; clinically meaningful threshold exceeded); \textbf{S} for microvascularity
&
Moderate (pilot)
&
\textbf{Industry-funded (GELITA\,AG / TENDOFORTE\textsuperscript{\textregistered}).} Flag.\\

\addlinespace

Khatri et al.\ (2021).
\textit{Amino Acids} 53:1647--1663. \cite{KNC+21}
&
Systematic review of 15 RCTs; healthy adults and clinical populations.
&
Collagen peptides 5--15\,g/day $+$ exercise, 3--12\,weeks.
&
Collagen synthesis markers, joint pain, body composition.
&
Pooled conclusion: collagen peptides $+$ exercise improve joint pain/function and increase collagen synthesis markers; no consistent additional effect on muscle protein synthesis vs.\ higher-quality proteins. Heterogeneous study designs; pooled ES not calculated.
&
Mixed: \textbf{S} (synthesis markers); \textbf{F} (joint pain/function outcomes in clinical populations)
&
Moderate (systematic review; heterogeneous risk of bias)
&
Some constituent trials industry-funded; authors note this.\\

\addlinespace

Buchalski et al.\ (2026).
\textit{J Funct Morphol Kinesiol} 11(1):130. \cite{BJAL26}%
\footnote{Cited by two independent sources; DOI should be verified via resolver before submission.}
&
Systematic review; 8 RCTs; n=257 (246M:11F); age 18--52\,yr; 3--15\,wk training.
&
Collagen peptides 5--30\,g/day $+$ ${\geq}50$\,mg vitamin~C $+$ high-load training.
&
Tendon cross-sectional area (CSA); stiffness; Young's modulus.
&
3/4 studies: CSA $\uparrow$ significantly vs.\ placebo (collagen: $+11\,\%$ vs.\ $+4.7\,\%$). 2/2 high-dose studies: stiffness and modulus $\uparrow$ significantly. No additive effect on muscular strength. GRADE A for structural outcomes.
&
\textbf{S-structural} (tendon CSA, stiffness, modulus via imaging/mechanics; not injury incidence or climbing performance; large tendons only)
&
Moderate (GRADE A for structural)
&
Not declared.\\

\addlinespace

Nulty et al.\ (2024).
\textit{Am J Physiol Endocrinol Metab}. \cite{NTD+24}%
\footnote{Single-source claim. Requires independent verification.}
&
Randomised crossover; n=8 resistance-trained men, age $49{\pm}8$\,yr.
&
0, 15, or 30\,g vitamin-C-enriched hydrolysed collagen 1\,h pre-resistance exercise.
&
P1NP AUC 0--6\,h.
&
0\,g: $96{\pm}23$\,\textmu g/mL${\times}$h; 15\,g: $134{\pm}23$ ($+40\,\%$, $p{=}0.024$); 30\,g: $169{\pm}28$ ($+76\,\%$ vs.\ 0\,g, $p{<}0.001$; $+26\,\%$ vs.\ 15\,g, $p{=}0.033$).
&
\textbf{S} (biomarker P1NP only; no structural or functional endpoint)
&
Low (small n; middle-aged; surrogate)
&
Not declared.\\

\addlinespace

Sas-Nowosielski et al.\ (2021).
\textit{Int J Environ Res Public Health} 18(10):5370. \cite{SNWK21}
&
Double-blind placebo-controlled; n=15 elite climbers (13M:2F, 1 excluded after injury); 4-week intervention.
&
$\beta$-alanine 4.0\,g/day vs.\ placebo; divided doses.
&
Campus-board intermittent reaches; easy traverse moves-to-failure and time-to-failure.
&
Campus reaches: $d{=}1.55$, $p{=}0.002$; easy traverse moves $+51\,\%$ ($d{=}1.73$); time-to-failure $+32$\,s ($d{=}1.44$, $p{=}0.044$). Hard traverse: no significant benefit.
&
Mixed: campus and traverse tasks are \textbf{S-functional} (climbing-specific but not grade or route completion); no on-route performance endpoint
&
Moderate (climbing-specific; small n; multiple endpoints)
&
None declared.\\

\addlinespace

Sandford et al.\ (2018).
\textit{BMJ Open Sport Exerc Med} 4(1):e000414. \cite{SSWL18}
&
Multicentre double-blind RCT; n=73 rotator-cuff-related shoulder pain; 8\,weeks + exercise.
&
9 capsules/day MaxEPA (1.53\,g EPA $+$ 1.04\,g DHA) vs.\ oleic acid placebo.
&
Oxford Shoulder Score (OSS) at 2 months; SPADI at 3 months.
&
OSS: mean difference $-0.1$ (95\,\%\,CI $-2.6$ to $2.5$; $p{=}0.95$) --- no effect. SPADI: $-8.3$ (95\,\%\,CI $-15.6$ to $-0.94$; $p{=}0.03$) at 3 months favoring omega-3.
&
\textbf{F} (validated pain/disability scores; shoulder not finger; null primary result)
&
Good (adequate n; multi-centre)
&
None declared.\\

\addlinespace

Grgic et al.\ (2022).
\textit{Clin Nutr ESPEN} 47:89--95. \cite{Grg22}
&
Systematic review and meta-analysis; 16 studies; n=353 (34 female).
&
Caffeine 1--7\,mg/kg vs.\ placebo; acute doses.
&
Isometric handgrip strength.
&
$d{=}0.17$ (95\,\%\,CI 0.10--0.23; $p{<}0.001$) overall. Subgroup 1--3\,mg/kg: $d{=}0.20$ (95\,\%\,CI 0.10--0.30; $p{<}0.001$). Small but significant ergogenic effect.
&
\textbf{S} (dynamometric grip strength; no climbing performance outcome; general athletic population)
&
Moderate (some risk of bias; male-predominant)
&
None declared.\\

\addlinespace

Ruiz-L\'opez et al.\ (2026).
\textit{Nutrients} 18(2):284. \cite{RLMAMR+26}
&
Triple-blind randomised crossover RCT; n=13 male climbers.
&
3\,mg/kg caffeine vs.\ placebo; administered 60\,min pre-testing.
&
Pull-up 1RM and power; pull-up endurance; dead-hang time; maximal dead-hang strength; RFD.
&
All outcomes non-significant: pull-up reps $+3.3\,\%$ (95\,\%\,CI $-3.30$ to $4.37$; $g{=}0.061$); maximal grip $+2.4\,\%$ ($p{=}0.060$; $g{=}0.120$); dead-hang $-7.4\,\%$ ($p{=}0.162$). Effect sizes trivial--small.
&
\textbf{S} (pull-up and dead-hang metrics are climbing-adjacent dynamometric surrogates; no on-wall or grade outcome)
&
Moderate (well-controlled; first climbing-specific caffeine RCT; underpowered)
&
None declared.\\

\end{longtable}

\noindent\textit{Note: No study in Table~\ref{tab:suppstudy} used a climbing-specific primary functional outcome (redpoint grade, on-wall endurance, pulley injury incidence). All climbing-specific evidence tiers are extrapolations.}

%---------------------------------------------------------------------
\subsection{Hydrolyzed Collagen Peptides}
\label{subsec:collagen}

\subsubsection{Mechanism}

Type~I collagen constitutes 65--80\,\% of tendon and ligament dry weight, and is the principal structural protein of finger flexor tendons, annular pulleys (A2, A4), and skin dermis --- all high-stress tissues in rock climbing. Collagen triple-helix formation requires hydroxylation of proline and lysine residues by prolyl-4-hydroxylase and lysyl hydroxylase, both of which are ascorbic-acid-dependent dioxygenases; without adequate vitamin~C, hydroxylation is rate-limited and secreted collagen is thermally unstable and subject to intracellular degradation. \textbf{[Evidence-backed: established biochemistry; Myllyharju, 2003, \textit{Matrix Biol}]}

Hydrolyzed collagen peptides (enzymatically processed to ${\approx}2$--5\,kDa; lower molecular weight than thermal gelatin) are absorbed intact from the gut as Pro-Hyp and Hyp-Gly dipeptides, which accumulate in connective tissue and may act as substrates for \textit{de novo} collagen synthesis and/or as signaling molecules upregulating tenocyte collagen gene expression. \textbf{[Bioplausible: mechanism established in vitro (Shigemura et al., 2009, \textit{J Agric Food Chem}); direct in-vivo mechanistic confirmation in human tendon tissue is lacking]}

Intermittent mechanical loading of tendons during exercise increases local blood flow and tenocyte metabolic activity, creating a window during which circulating amino-acid precursors are preferentially incorporated. The Shaw (2017) protocol exploited this by timing collagen ingestion 1\,h pre-exercise to synchronise peak plasma amino-acid elevation with exercise-induced hyperaemia. \textbf{[Bioplausible: inferred from P1NP kinetics; direct tissue-incorporation measurements absent]}

\subsubsection{Shaw et al.\ (2017) --- Critical Appraisal}

Shaw et al.\ (\textit{Am J Clin Nutr} 105:136--143; \cite{SLBR+17}) used a randomised, double-blind, crossover design in n=8 healthy males (age $21{\pm}0.5$\,yr; no tendon pathology). Three conditions --- 5\,g gelatin\,$+$\,48\,mg vitamin~C, 15\,g gelatin\,$+$\,48\,mg vitamin~C, and placebo --- were each administered as a single dose 1\,h before 6\,min rope-skipping, repeated 3$\times$/day for 3 days, with a 7-day washout.

\textbf{Primary outcomes}: (1) Serum P1NP --- a systemic surrogate biomarker of type~I collagen synthesis rate, \textit{not} a direct measure of tendon or pulley collagen content; (2) mechanical properties (stiffness, collagen content) of in-vitro engineered ligament constructs incubated with participant serum --- a cell-culture functional proxy, not a measurement in the participants' own tendons.

\textbf{Results}: The 15\,g condition produced ${\approx}2{\times}$ greater P1NP AUC vs.\ placebo over 4\,h post-exercise ($p{<}0.05$). Engineered ligament stiffness and collagen content were significantly greater with post-15\,g serum vs.\ post-placebo serum. The 5\,g condition was intermediate. Formal Cohen's $d$ and 95\,\%\,CI were not reported; the 2$\times$ P1NP estimate is a within-person comparison from n=8 with low precision.

\textbf{Critical limitations}:
\begin{enumerate}[noitemsep]
  \item n=8, male-only: underpowered; female hormonal effects on collagen turnover (estrogen upregulates collagen synthesis) unaddressed.
  \item Surrogate endpoint: P1NP reflects systemic whole-body collagen synthesis and does not confirm preferential incorporation into finger pulleys or flexor tendons. No tendon biopsy data exist.
  \item Healthy tendons only: tendinopathic tissue has an altered tenocyte phenotype and remodeling dynamics; results cannot be directly extrapolated to injured A2/A4 pulleys.
  \item Acute single-dose design: no chronic dosing, MRI, ultrasound, or functional load-to-failure outcome was assessed.
  \item Gelatin $\ne$ hydrolyzed collagen peptides: gelatin (thermally denatured) has higher molecular weight fragments than enzymatically hydrolyzed collagen (${\approx}2$--5\,kDa). Lis \& Baar (2019) found comparable amino-acid responses when both were vitamin-C-enriched, suggesting functional equivalence under those conditions.
  \item No climbing-specific outcome: rope-skipping loads Achilles/patellar tendons, not finger flexor pulleys.
  \item Industry conflict: the Baar lab has received collagen-industry support; a co-author holds patents on engineered ligament technology. Not disqualifying, but warrants bias assessment.
\end{enumerate}

\begin{evidencebox}{Collagen Pre-Exercise: P1NP Surrogate Benefit (Shaw 2017) · \evidencebacked · ESS 4}
\textbf{Evidence-backed (surrogate biomarker only):} 15\,g gelatin $+$ 48\,mg vitamin~C taken 1\,h before intermittent loading acutely elevates systemic P1NP in healthy young males (${\approx}2{\times}$; $p{<}0.05$) and improves in-vitro engineered ligament mechanics. \textit{Endpoint: surrogate (P1NP biomarker; ex-vivo ligament stiffness). No direct human tendon structural outcome, no injury incidence, no return-to-climbing, no climbing performance endpoint. This does not constitute evidence that climbers heal pulleys faster, increase tendon stiffness, or reduce injury incidence.}
\end{evidencebox}

\subsubsection{Lis \& Baar (2019) Confirmation}

Lis \& Baar (\textit{IJSNEM} 29:526--531; \cite{LB19}): n=10 males, crossover; all active collagen conditions (vitamin-C-enriched gelatin, hydrolyzed collagen, gummy) increased circulating amino acids (glycine, proline, hydroxyproline) similarly; P1NP showed a ${\approx}20\,\%$ trend from baseline that did not reach statistical significance due to high variability. The study confirms the amino-acid-kinetic window (peak ${\approx}1$\,h post-ingestion) and suggests functional equivalence between gelatin and hydrolyzed collagen forms when vitamin~C is co-supplemented. It does not establish a statistically reliable P1NP effect at 15\,g hydrolyzed collagen alone; the non-significance limits claims of confirmed replication.

\subsubsection{Praet et al.\ (2019): Clinical Tendinopathy Data}

Praet et al.\ (\textit{Nutrients} 11:76; \cite{PPW+19}) tested 5\,g/day specific collagen peptides (TENDOFORTE\textsuperscript{\textregistered}) $+$ calf-strengthening vs.\ exercise-only in n=20 chronic Achilles tendinopathy patients (pilot, double-blind, crossover RCT, 6 months). VISA-A improved by $+12.6$ (95\,\%\,CI 9.7--15.5) in the collagen-first group vs.\ $+5.3$ (95\,\%\,CI 2.3--8.3) in the placebo-first group; the between-group difference exceeds the MCID of ${\approx}6.5$ points. Tendon microvascularity decreased similarly in both groups.

\textbf{Limitations}: Pilot scale (n=20); crossover design with potential carry-over effects; single commercial peptide brand (TENDOFORTE\textsuperscript{\textregistered}); \textbf{funded by GELITA AG} (industry-funded: flag); Achilles tendon loading differs substantially from finger flexor pulley loading in climbing; extrapolation to A2/A4 pulleys requires a tissue-type and loading-pattern leap.

\begin{bioplausbox}{Collagen Peptides + Loading for Connective Tissue Recovery in Climbers · \bioplausible · ESS 2}
\textbf{Bioplausible} for collagen peptides $+$ structured loading supporting connective tissue recovery in climbers, drawing on Praet et al.\ (2019). The closest available analog involves a different tendon, a lower dose, a different loading protocol, and an industry-funded trial. No direct evidence exists for climbing-specific pulley healing.
\end{bioplausbox}

\subsubsection{Dose}

Shaw (2017) demonstrated a dose-response between 5\,g and 15\,g for P1NP elevation; 15\,g produced the strongest signal. Lis \& Baar (2019) used 15\,g for all active conditions. Khatri et al.\ (2021; systematic review, 15 RCTs) identified 5--15\,g/day as the studied range, with 15\,g consistently associated with the highest collagen-synthesis marker responses. Praet et al.\ (2019) showed clinical efficacy at 5\,g/day in a chronic tendinopathy context with structured loading, suggesting lower doses may be sufficient for chronic daily protocols. No RCT has tested doses $>$15\,g in young healthy athletes; evidence from one single-source study (Nulty et al., 2024)\footnote{Single-source claim --- requires independent verification.} in middle-aged men suggests 30\,g produces $+26\,\%$ greater P1NP AUC than 15\,g ($p{=}0.033$), but this has not been replicated.

\textbf{[Evidence-backed:]} 15\,g is the best-supported acute pre-exercise dose for collagen synthesis biomarker elevation. No evidence that doses $>$15\,g provide additional clinical benefit in young athletes.

\subsubsection{Timing Window}

Shaw (2017) and Lis \& Baar (2019) both timed collagen ingestion 1\,h before exercise to align the plasma amino-acid peak (${\approx}60$\,min post-ingestion for 15\,g) with exercise-induced tendon hyperaemia. Gastric emptying variability (affected by co-ingested food, hydration, and sympathetic tone) could shift the peak by ${\pm}15$--30\,min. Taking collagen 45--75\,min pre-session adequately approximates the Shaw timing window. Immediate pre-exercise ingestion ($<$15\,min) likely misaligns the amino-acid peak with early high-strain loading.

\textbf{[Bioplausible:]} 1\,h pre-exercise is mechanistically optimal; direct evidence that the 1\,h window is superior to, e.g., 45\,min or 90\,min, does not exist from RCT comparisons.

\subsubsection{The Vitamin~C Co-Ingestion Gap}
\label{subsubsec:vitc}

\begin{warningbox}{\textbf{Protocol Gap: Vitamin\textasciitilde{}C Absent from Collagen Shake}}
The user takes 15\,g hydrolyzed collagen peptides pre-session \textbf{without} vitamin~C co-ingestion. Shaw et al.\ (2017) and Lis \& Baar (2019) both used vitamin-C-enriched preparations (48\,mg ascorbic acid per dose). Without this co-factor, the evidence-backed biomarker benefit \textbf{does not cleanly transfer}.
\end{warningbox}

\textbf{Biochemistry}: Prolyl-4-hydroxylase requires ascorbic acid as an obligatory electron donor; in its absence the hydroxylation cycle stalls, and hydroxylation efficiency falls in proportion to ascorbate depletion. Under frank deficiency ($<$10\,mg/day), collagen triple-helix formation fails entirely (scurvy). At suboptimal non-scorbutic intakes (plasma ascorbate $<$23\,$\mu$mol/L), hydroxylation efficiency is proportionally reduced. \textbf{[Evidence-backed: established enzymology]}

\textbf{Shaw protocol context}: 48\,mg vitamin~C was co-ingested per dose. The US RDA for vitamin~C is 90\,mg/day for men and 75\,mg/day for women; 48\,mg is sub-RDA but was not intended as a pharmacological dose --- its purpose was to standardise local ascorbate availability during the post-ingestion window without confounding the collagen signal.

\textbf{Scenario analysis} for the user's protocol:

\begin{center}
\begin{tabular}{>{\raggedright\arraybackslash}p{3.2cm}
                >{\raggedright\arraybackslash}p{3.8cm}
                >{\raggedright\arraybackslash}p{5.5cm}}
\toprule
\textbf{Vitamin~C status} & \textbf{Plasma ascorbate} & \textbf{Expected collagen synthesis response} \\
\midrule
Adequate habitual diet (${\geq}5$ fruit/veg servings/day; ${\geq}100$\,mg/day dietary C) & ${\geq}50$\,$\mu$mol/L (replete) & Prolyl hydroxylase likely not substrate-limited; response may approximate Shaw protocol. Co-ingestion provides marginal additional ascorbate. \\
\addlinespace
Suboptimal intake (common in athletes in caloric restriction; $<$50\,mg/day dietary C) & ${<}23$\,$\mu$mol/L (sub-optimal) & Substrate (Pro, Gly, Hyp) elevated but hydroxylation rate-limited; collagen synthesis efficiency impaired; evidence-backed benefit does not apply. \\
\bottomrule
\end{tabular}
\end{center}

\textbf{Critical gap}: No RCT has directly compared 15\,g collagen \textit{with} vs.\ \textit{without} co-ingested vitamin~C under controlled dietary ascorbate conditions. The claim that omitting vitamin~C definitively abolishes benefit is \textbf{[Bioplausible]}, not experimentally proven. However, the claim that the user's protocol replicates the Shaw design is \textbf{false}: the designs differ in a mechanistically important cofactor. The uncertainty is asymmetric --- co-ingesting 50--100\,mg vitamin~C has no downside (cost $<$\$0.01/day; no risk at this dose) and eliminates the bottleneck. \textbf{Absence of co-ingestion is a correctable gap with no counterargument.}

\textbf{[Bioplausible]:} Adding ${\approx}50$\,mg ascorbic acid to the collagen shake is the single highest-priority protocol modification for the user. Whether it is strictly necessary depends on habitual dietary vitamin~C status, which has not been measured.

\subsubsection{Claim: Collagen for Climbing Skin Recovery}

\textbf{Claim}: Oral hydrolyzed collagen peptides accelerate healing of flappers (partial-thickness skin avulsions), skin tears, and chronic fingertip skin thinning in climbers.

\textbf{Evidence}: Puškarić et al.\ (2024; \textit{Nutrients}, \cite{FLAGA+24})\footnote{\toverify: only two sources cite a skin-specific RCT; the DOI cited by other sources resolves to Fernández-Lázaro 2024 omega-3 review. The skin RCT is a single-source claim requiring independent DOI verification.} examined collagen (5\,g $+$ 80\,mg vitamin~C/day $\times$ 16\,wk) in aging women (n=87), showing improved dermis density and texture vs.\ placebo. This population (cosmetic aging) and outcome (dermal density) differ mechanistically from acute traumatic skin avulsion and callus reformation in climbers.

No peer-reviewed RCT has tested oral collagen peptides for climbing-specific skin injuries.

\begin{folklorebox}{Collagen Accelerates Healing of Climbing Flappers and Split Tips · \folkloreverified · ESS 1}
\textbf{Folklore:} ``Collagen accelerates healing of climbing flappers and split tips.'' Source: climbing community forums (r/climbharder), coach blogs (Hooper's Beta, Lattice Training social media), supplement marketing. No peer-reviewed climbing-specific or acute wound-healing evidence exists for this claim. Bioplausible substrate mechanism (collagen provision $\to$ dermal fibroblast activity) does not constitute evidence.
\end{folklorebox}

\subsubsection{Claim: Collagen for Pulley and Tendon Remodeling Under Active Climbing Load}

\textbf{Mechanism} \textbf{[Bioplausible]}: Tendons and annular pulleys are ${\approx}70\,\%$ type~I collagen by dry weight. During loading, tenocytes upregulate TGF-$\beta$1 and IGF-1 signaling, increasing collagen turnover. Net anabolic balance requires adequate substrate (proline, glycine, hydroxyproline) and co-factors (vitamin~C). Provision of 15\,g hydrolyzed collagen elevates circulating Gly ($+376$\,$\mu$mol/L), Pro ($+162$\,$\mu$mol/L), and Hyp ($+105$\,$\mu$mol/L) at ${\approx}1$\,h post-ingestion (Shaw, 2017); active climbing concurrently increases tendon perfusion. The anabolic-window model is mechanistically coherent. \textbf{[Bioplausible: mechanism plausible; no climbing RCT measures pulley or flexor tendon collagen content change attributable to oral supplementation]}

Buchalski et al.\ (2026; systematic review) report that collagen 15--30\,g/day $+$ vitamin~C $+$ high-load training (${\geq}70\,\%$ 1RM resistance) increases tendon CSA ($+11\,\%$ vs.\ $+4.7\,\%$ placebo) and stiffness/modulus in 3/4 studies; GRADE A for structural outcomes.\footnote{Two independent sources cite Buchalski et al.\ 2026. Source metadata requires resolver verification before submission.} This is the strongest structural-outcome evidence available, but it involves large tendons under resistance-training load, not small-diameter finger pulleys under climbing-specific isometric crimp load.

\begin{bioplausbox}{Collagen for Pulley and Tendon Remodeling Under Climbing Load · \bioplausible · ESS 2}
\textbf{Bioplausible:} Collagen peptides $+$ vitamin~C $+$ structured loading may support type~I collagen anabolism in finger flexor tendons and A2/A4 pulleys. The Praet (2019) Achilles tendinopathy data and Buchalski (2026) structural data are the closest available analogs. No climbing-specific RCT exists for pulley healing, injury prevention, or return-to-climbing outcomes.
\end{bioplausbox}

\subsubsection{User Protocol Assessment and Recommended Correction}

\textbf{Current protocol}: 15\,g hydrolyzed collagen peptides pre-climbing session, no vitamin~C, unspecified exact timing.

\begin{center}
\begin{tabular}{>{\raggedright\arraybackslash}p{3.5cm}
                >{\raggedright\arraybackslash}p{2.5cm}
                >{\raggedright\arraybackslash}p{3.2cm}
                >{\raggedright\arraybackslash}p{3.5cm}}
\toprule
\textbf{Parameter} & \textbf{User protocol} & \textbf{Evidence-backed protocol} & \textbf{Gap} \\
\midrule
Dose & 15\,g \checkmark & 15\,g (Shaw/Lis/Khatri) & None \\
\addlinespace
Vitamin~C & \textbf{None} & ${\geq}48$--50\,mg co-ingested & \textbf{Critical gap} \\
\addlinespace
Timing & ``Before climbing'' (unspecified) & 45--75\,min pre-session & Clarify; risk if $<$15\,min \\
\addlinespace
Frequency & Session days only & Daily (Shaw: 3$\times$/day $\times$3\,d; Baar 2017 recommendation: daily) & Uncertain gap; rest-day dosing untested \\
\addlinespace
Loading stimulus & Rock climbing & Intermittent loading ${\geq}6$\,min at moderate--high intensity & Climbing provides adequate stimulus \checkmark \\
\addlinespace
Monitoring & Unknown & Periodic grip/pain VAS assessment & No monitoring \\
\bottomrule
\end{tabular}
\end{center}

\textbf{Minimum corrective action}: Co-ingest 50--100\,mg ascorbic acid (e.g., 100\,mg tablet or a small glass of citrus juice) with the 15\,g collagen shake, taken 45--60\,min before session start. This aligns the protocol with Shaw et al.\ (2017) and eliminates the mechanistic bottleneck at prolyl hydroxylase.

%---------------------------------------------------------------------
\subsection{Creatine Monohydrate}
\label{subsec:creatine}

\textbf{Mechanism}: Creatine supplementation saturates intramuscular phosphocreatine (PCr) stores, accelerating ATP resynthesis during high-intensity efforts of $\leq$15\,s (PCr-dependent energy system). Secondary: creatine may augment muscle hypertrophy via cell volumisation and mTOR pathway activation. Relevant to climbing: maximal boulder moves, contact-strength efforts, campus board power, and repeated attempts with brief rest intervals. \textbf{[Evidence-backed: well-replicated mechanism]}

\textbf{Evidence}: Multiple meta-analyses document creatine's efficacy for muscular strength and power in resistance-trained individuals. No climbing-specific RCT has measured creatine effects on boulder problem performance, maximum hang time, or campus board output in trained climbers. \textbf{[Evidence-backed for general strength/power; Bioplausible for climbing]}

\textbf{Dose}: No loading phase required for saturation if 3--5\,g/day maintained for $\geq$28 days. If rapid saturation is desired: 20\,g/day $\div$ 4 doses $\times$ 5--7 days, then 3--5\,g/day maintenance. Creatine monohydrate is the only form with robust evidence; effervescent, ethyl ester, and buffered variants show no superiority.

\textbf{Climbing-specific caveat}: Body mass increases of ${\approx}$0.5--2\,kg from intracellular water retention may impair strength-to-weight ratio and thus on-wall performance, particularly in weight-sensitive disciplines (lead climbing, slabs, long traverses). The net effect for power-oriented bouldering disciplines is likely positive, but this has not been tested in a controlled climbing trial.

\begin{evidencebox}{Creatine: Strength and Power in Resistance-Trained Athletes · \evidencebacked · ESS 4}
\textbf{Evidence-backed} for general strength and power in resistance-trained athletes (multiple meta-analyses). \textbf{Bioplausible} for climbing performance improvement; body-mass caveat is real and applies specifically to climbing.
\end{evidencebox}

\begin{toverifybox}{Creatine Body-Mass Penalty in Climbing Performance · \toverify · ESS 1}
\textbf{\toverify:} Community claims that ``creatine makes you too heavy to climb'' are anecdotal and not supported by any controlled climbing-performance study. The magnitude of mass gain (0.5--2\,kg) and its effect on climbing-specific metrics remain empirically untested. Only two sources identified (r/climbharder; commercial supplement blogs); requires $\geq$3 independent climbing-community sources for \folkloreverified status.
\end{toverifybox}

%---------------------------------------------------------------------
\subsection{Omega-3 Fatty Acids (EPA/DHA)}
\label{subsec:omega3}

\textbf{Mechanism}: EPA and DHA compete with arachidonic acid for COX-2 and LOX enzymatic pathways, shifting eicosanoid and leukotriene profiles toward less pro-inflammatory mediators. Potential relevance to climbing: modulating tendon/ligament inflammation in overuse pathology (pulley tendinopathy, medial epicondylopathy), reducing training-induced DOMS, and supporting synovial joint health in finger DIP/PIP joints. \textbf{[Evidence-backed: biochemical mechanism; Calder, 2017, \textit{Eur J Clin Nutr}]}

\textbf{Best tendon-specific RCT evidence}: Sandford et al.\ (2018; \cite{SSWL18}), n=73 rotator-cuff-related shoulder pain, 2.57\,g/day EPA+DHA $\times$ 8 weeks $+$ exercise: no clinically meaningful OSS benefit at 2 months (mean difference $-0.1$, 95\,\%\,CI $-2.6$ to $2.5$; $p{=}0.95$); modest SPADI improvement at 3 months ($-8.3$, 95\,\%\,CI $-15.6$ to $-0.94$; $p{=}0.03$) of uncertain clinical significance. No direct benefit for structural tendon outcomes. \textbf{[Evidence-backed for systemic anti-inflammatory effect; Evidence-backed (null result) for tendinopathy structural/pain outcomes in this population]}

A Mendelian randomisation study found genetically higher omega-3 levels were associated with an $+18\,\%$ increase in Achilles tendinitis risk.\footnote{Single-source claim --- requires independent verification. Mendelian randomisation cannot establish causality and is subject to pleiotropy; interpret with caution.}

\textbf{Dose}: 2--3\,g/day combined EPA+DHA is the dose range used in most anti-inflammatory trials. No dose-response curve for tendon outcomes. Daily supplementation; timing is not acutely critical. Antiplatelet effect at $>$3\,g/day; anticoagulant interaction at high doses; fish oil peroxidation if stored improperly; algal sources available for vegetarians.

\begin{bioplausbox}{Omega-3 for Tendon Inflammation Modulation in Climbers · \bioplausible · ESS 3}
\textbf{Bioplausible} for modulation of tendon-related inflammation in climbers with overuse pathology. Current controlled evidence does not confirm clinically meaningful tendon-structural or injury-prevention benefit. The best available tendinopathy RCT (Sandford 2018, rotator cuff) shows no primary-outcome benefit. No climbing-specific RCT exists.
\end{bioplausbox}

\begin{toverifybox}{Omega-3 Strengthens Tendons or Prevents Pulley Injuries in Climbers · \toverify · ESS 0}
\textbf{\toverify:} Claims that omega-3 ``strengthens tendons'' or ``prevents pulley injuries'' in climbers. Only two sources identified (physiotherapy blogs; sports nutrition marketing); requires $\geq$3 independent climbing-community sources. Not supported by current controlled evidence (Sandford 2018 null result).
\end{toverifybox}

%---------------------------------------------------------------------
\subsection{Magnesium}
\label{subsec:magnesium}

\textbf{Mechanism}: Mg$^{2+}$ is cofactor for $>$300 enzymatic reactions, including Mg-ATP complex formation, muscle relaxation (antagonises Ca$^{2+}$ at the sarcoplasmic reticulum), and sleep architecture modulation (GABA-receptor agonism). Relevant to climbing: neuromuscular recovery between moves, sleep quality for training adaptation, and potentially cramp prevention. \textbf{[Evidence-backed: physiological role]}

\textbf{Evidence}: Magnesium supplementation corrects documented deficiency and may improve performance in deficient individuals; effects in magnesium-replete athletes are small or absent. Cochrane review and a JAMA RCT (n=94) found magnesium oxide no better than placebo for nocturnal leg cramps ($\Delta -0.38$ cramps/week in both groups).\footnote{Single-source claim for the JAMA cramp RCT --- requires independent verification.} For sleep, a specific magnesium bisglycinate (250\,mg elemental/day) trial reported Insomnia Severity Index improvement ($d{\approx}0.2$, CI reported) at 4 weeks in poor sleepers.\footnote{Single-source claim --- requires independent verification.} No climbing-specific RCT exists for magnesium effects on forearm pump, grip endurance, or training recovery. \textbf{[Bioplausible for climbing performance benefit in athletes with suboptimal status]}

Climbers in caloric restriction, heavy sweating, and high training loads are at elevated risk of suboptimal (not frank deficiency) magnesium status. Serum magnesium is an insensitive biomarker; erythrocyte or ionised magnesium testing is more informative.

\textbf{Dose}: 200--400\,mg elemental magnesium/day (as glycinate, malate, or citrate; superior bioavailability vs.\ oxide). Evening dosing to support sleep benefit. Diarrhoea common at $>$400\,mg oxide.

\begin{bioplausbox}{Magnesium for Sleep and Neuromuscular Recovery in Suboptimal-Status Climbers · \bioplausible · ESS 2}
\textbf{Bioplausible} for sleep-mediated recovery benefit in climbers with suboptimal magnesium status. \textbf{Evidence-backed} for correcting documented deficiency. \textbf{Evidence-backed (null)} for cramp prevention in general populations.
\end{bioplausbox}

\begin{toverifybox}{Magnesium Prevents Forearm Pump or Finger Cramps in Climbers · \toverify · ESS 0}
\textbf{\toverify:} Claims that magnesium ``prevents forearm pump'' or ``stops finger cramps on routes.'' Only two sources identified (coaching forums; gym culture); requires $\geq$3 independent climbing-community sources. No targeted trial supports this in euvolaemic, well-fed climbers.
\end{toverifybox}

%---------------------------------------------------------------------
\subsection{Beta-Alanine}
\label{subsec:betaalanine}

\textbf{Mechanism}: $\beta$-alanine is the rate-limiting precursor for muscle carnosine synthesis. Carnosine ($\beta$-alanyl-L-histidine; pKa ${\approx}$6.83) acts as an intramuscular H$^+$ buffer, attenuating acidosis during high-intensity glycolytic work. Forearm pump in rock climbing --- the accumulating burning sensation from sustained gripping on crux sections and power-endurance routes --- is mechanistically attributable to intramuscular H$^+$ accumulation from glycolysis; carnosine's buffering capacity targets precisely this mechanism. \textbf{[Bioplausible for climbing; mechanism Evidence-backed in general contexts]}

\textbf{Evidence}: Hobson et al.\ (2012, \textit{Amino Acids}, meta-analysis, k=15 RCTs): $\beta$-alanine significantly increased exercise capacity for efforts of 60--240\,s (ES = 0.374, 95\,\%\,CI 0.140--0.607). Trexler et al.\ (2015, \textit{J Int Soc Sports Nutr}, meta-analysis, k=57 studies): mean performance improvement $+2.85\,\%$ (95\,\%\,CI 0.37--10.49); Hedges $g = 0.374$ for 60--240\,s efforts specifically. This duration range (1--4\,min) corresponds to sustained pump sections and power-endurance sequences on sport climbing routes. \textbf{[Evidence-backed for 60--240\,s high-intensity exercise]}

\textbf{Climbing-specific RCT}: Sas-Nowosielski et al.\ (2021; \cite{SNWK21}): $\beta$-alanine 4\,g/day $\times$ 4 weeks in n=15 elite climbers (double-blind, placebo-controlled). Campus-board reaches: $d{=}1.55$, $p{=}0.002$; easy traverse moves: $+51\,\%$, $d{=}1.73$; easy traverse time-to-failure: $+32$\,s, $d{=}1.44$. Hard traverse: no significant benefit. The large effect sizes should be interpreted cautiously given the small sample (n=15, with 1 excluded) and multiple endpoints without correction for multiplicity. \textbf{[Evidence-backed for specific climbing-relevant endurance tasks; small n; moderate quality]}

\textbf{Dose}: 3.2--6.4\,g/day in divided doses (0.8--1.6\,g per dose, every 3--4\,h) to minimise paresthesia. Carnosine loading requires 4--12 weeks at therapeutic doses for ${\approx}$60--80\,\% increase in muscle carnosine content.

\textbf{Caveat}: Paresthesia (harmless tingling, primarily face/scalp/hands) is near-universal at single doses $>$800\,mg; divided or slow-release formulations mitigate this. Performance benefit is most pronounced for efforts of 60--240\,s; submaximal long-duration routes with low anaerobic contribution may see less benefit. No benefit for maximal strength or power output.

\begin{evidencebox}{Beta-Alanine: Climbing Endurance Task Benefit (Sas-Nowosielski 2021) · \evidencebacked · ESS 7}
\textbf{Evidence-backed} for 60--240\,s high-intensity exercise (multiple meta-analyses, Hedges $g = 0.374$) and for specific climbing-endurance tasks in elite climbers (campus board, easy traverse; Sas-Nowosielski 2021, $d = 1.44$--1.73, but n=15 only). Dose: 3.2--6.4\,g/day divided $\times \geq$4 weeks. \textit{Endpoint: campus-board and traverse tasks are climbing-specific surrogate-functional hybrids; no redpoint grade, route completion rate, or injury incidence endpoint. Hard traverse showed no benefit.}
\end{evidencebox}

\begin{toverifybox}{Beta-Alanine Eliminates Forearm Pump in Climbing · \toverify · ESS 0}
\textbf{\toverify:} Claims that $\beta$-alanine ``eliminates pump'' in climbing. Only two sources identified (climbing blogs; nutrition forums); requires $\geq$3 independent climbing-community sources. Current data show improvement in specific tasks, not elimination of pump; hard traverse performance showed no benefit in the only climbing RCT (Sas-Nowosielski 2021).
\end{toverifybox}

%---------------------------------------------------------------------
\subsection{Caffeine}
\label{subsec:caffeine}

\textbf{Mechanism}: Non-selective adenosine receptor antagonist (A1, A2A). Reduces perceived exertion, enhances neuromuscular activation, increases pain tolerance, potentiates motor unit recruitment, and accelerates calcium release from the sarcoplasmic reticulum. Relevant to climbing: sustained isometric grip force, delayed onset of perceived forearm pump, acute power output on difficult moves. \textbf{[Evidence-backed]}

\textbf{Best general evidence}: Grgic et al.\ (2022, \textit{Clin Nutr ESPEN} 47:89--95; \cite{Grg22}): meta-analysis, k=16 studies, n=353; isometric handgrip strength: $d{=}0.17$ (95\,\%\,CI 0.10--0.23; $p{<}0.001$). Subgroup 1--3\,mg/kg: $d{=}0.20$ (95\,\%\,CI 0.10--0.30). Small but statistically significant ergogenic effect on a measurement directly analogous to climbing grip. \textbf{[Evidence-backed for isometric grip strength]}

Grgic et al.\ (2020, \textit{Br J Sports Med} 54:681--688; \cite{WTNZ18}): k=34 studies; muscle strength ES = 0.20 (95\,\%\,CI 0.13--0.26); muscular endurance ES = 0.41 (95\,\%\,CI 0.24--0.58). \textbf{[Evidence-backed for general muscular endurance]}

\textbf{Climbing-specific RCT}: Ruiz-L\'opez et al.\ (2026; \cite{RLMAMR+26}): triple-blind randomised crossover, n=13 male climbers, 3\,mg/kg caffeine vs.\ placebo, 60\,min pre-testing. Pull-up reps: $+3.3\,\%$ (95\,\%\,CI $-3.30$ to $4.37$; $g{=}0.061$); maximal grip: $+2.4\,\%$ ($p{=}0.060$; $g{=}0.120$); dead-hang: $-7.4\,\%$ ($p{=}0.162$). All outcomes non-significant; effect sizes trivial--small. Subjective energy and alertness increased ($p{<}0.05$). This is the only published climbing-specific caffeine RCT and is underpowered (n=13); it does not rule out a small real effect at higher doses. \textbf{[Evidence-backed: 3\,mg/kg does not meaningfully improve proxy climbing performance at this sample size]}

\textbf{Dose}: 3--6\,mg/kg body mass, 45--60\,min pre-session. Lower doses (3\,mg/kg) sufficient for endurance and perceived-exertion benefit with fewer side effects; higher doses (5--6\,mg/kg) may increase maximal strength/power but risk anxiety, tachycardia, and sleep disruption. Habitual consumers show attenuated response; 48--96\,h abstinence before target sessions partially restores the effect.

\textbf{Climbing-specific caveat}: Fear and anxiety are performance factors in climbing, particularly on serious routes and in competition. Caffeine increases sympathetic arousal and may exacerbate anxiety-driven performance decrements, tremor, and hesitation on high-commitment moves. Sleep disruption at doses $>$3\,mg/kg within 6\,h of sleep onset reduces recovery quality.

\begin{evidencebox}{Caffeine: Isometric Grip Strength and Endurance Benefit · \evidencebacked · ESS 4}
\textbf{Evidence-backed} for isometric handgrip strength ($d{=}0.17$, Grgic 2022) and general muscular endurance ($d{=}0.41$, Grgic 2020). \textbf{Evidence-backed (null at 3\,mg/kg, n=13)} for climbing-specific pull-up/grip metrics (Ruiz-L\'opez 2026). Dose 3--6\,mg/kg 45--60\,min pre-session; individual tolerance and habituation are critical modifiers. \textit{Endpoint: dynamometric surrogate (grip strength, endurance metrics); no redpoint grade, route completion, or competition ranking outcome in any trial. The only climbing-specific RCT (Ruiz-L\'opez 2026) showed null results on all proxy climbing metrics.}
\end{evidencebox}

%---------------------------------------------------------------------
\subsection{What the Literature Cannot Currently Resolve}
\label{subsec:suppgaps}

\begin{enumerate}[noitemsep]

  \item \textbf{Collagen $+$ vitamin~C and pulley structure in climbers}: No study has measured A2/A4 pulley or finger flexor tendon CSA, stiffness, or load-to-failure as a function of oral collagen supplementation using high-resolution ultrasound, MRI, or biopsy in any climbing population. All positive structural data (Buchalski 2026) involve large tendons under resistance-training loads; small-diameter, high-strain pulley responses are unknown.

  \item \textbf{Minimum effective vitamin~C co-dose}: Shaw (2017) used 48\,mg; Lis (2019) used enriched preparations. No factorial RCT has compared collagen $\pm$ vitamin~C under controlled dietary ascorbate conditions, so the threshold above which co-ingestion is redundant for athletes with replete diets remains undefined.

  \item \textbf{Dose-response above 15\,g}: No RCT has compared 15\,g vs.\ 20\,g vs.\ 30\,g collagen with functional tendon outcomes in young healthy athletes. The 30\,g data (Nulty 2024, single source, middle-aged men) suggest an additional P1NP increment but have not been independently replicated.

  \item \textbf{Collagen and climbing skin repair}: No controlled trial has assessed time-to-heal for flappers, split tips, or chronic skin thinning in climbers receiving collagen supplementation vs.\ placebo under standardised wound-care conditions.

  \item \textbf{Beta-alanine and perceived pump in climbing}: The Sas-Nowosielski (2021) RCT showed campus board and easy traverse improvements but no hard traverse benefit. Whether carnosine loading reduces perceived pump and improves on-route performance in trained sport climbers (where vascular occlusion, not only metabolic acidosis, contributes to pump) is not established.

  \item \textbf{Caffeine at doses $>$3\,mg/kg in climbing}: The only climbing-specific RCT (Ruiz-L\'opez 2026) used 3\,mg/kg only. Whether 5--6\,mg/kg produces clinically meaningful improvement in actual climbing performance (redpoint grade, route completion percentage, competition ranking) without offsetting anxiety or tremor remains untested in climbing populations.

  \item \textbf{Omega-3 and finger tendon remodeling}: The available tendinopathy RCT (Sandford 2018, rotator cuff) showed no primary-outcome benefit. No RCT has examined EPA/DHA effects specifically on finger flexor tendons or pulleys, or on A2 pulley injury recurrence rate.

  \item \textbf{Interaction effects}: No study has examined whether collagen $+$ vitamin~C interacts with creatine, beta-alanine, or caffeine in terms of tendon adaptation, injury risk, or climbing-specific performance in a factorial or combination design.

\end{enumerate}

%---------------------------------------------------------------------
\section{Finger Injury Management}
\label{sec:injury}

\begin{warningbox}{Evidence Caveat: No Climbing-Specific Pulley Rehabilitation RCTs}
No randomized controlled trials (RCTs) have been conducted on climbing pulley-tear rehabilitation protocols, return-to-climb criteria, or taping versus no-taping for clinical outcomes. Climbing-specific intervention evidence exists for hangboard training and is referenced where applicable, but pulley-specific rehabilitation protocols are classified \textbf{Bioplausible} (mechanistically extrapolated from non-climbing tendinopathy RCTs or tissue biology) or \textbf{Folklore} (practitioner/community heuristic). Evidence-backed claims are restricted to diagnostic imaging accuracy, biomechanics studies, the Schneeberger \& Schweizer (2016) conservative splint series, and the Sjöman et al.\ (2023) injury-risk survey. This limitation must be front-loaded in any clinical application of these protocols.
\end{warningbox}

%---------------------------------------------------------------------
\subsection*{Summary Table}

{\footnotesize
\setlength{\tabcolsep}{4pt}
\renewcommand{\arraystretch}{1.2}
\begin{longtable}{>{\raggedright\arraybackslash}p{2.0cm}
                  >{\raggedright\arraybackslash}p{1.8cm}
                  >{\raggedright\arraybackslash}p{1.8cm}
                  >{\raggedright\arraybackslash}p{2.6cm}
                  >{\raggedright\arraybackslash}p{2.8cm}
                  >{\raggedright\arraybackslash}p{1.2cm}
                  >{\raggedright\arraybackslash}p{1.8cm}}
\caption{Peer-reviewed studies cited in this section. S = surrogate; F = functional.}\label{tab:injury_studies}\\
\toprule
\textbf{Authors} & \textbf{Design / $n$} & \textbf{Injury type} & \textbf{Primary outcome} & \textbf{Effect size / key stat} & \textbf{S/F} & \textbf{Quality} \\
\midrule
\endfirsthead
\toprule
\textbf{Authors} & \textbf{Design / $n$} & \textbf{Injury type} & \textbf{Primary outcome} & \textbf{Effect size / key stat} & \textbf{S/F} & \textbf{Quality} \\
\midrule
\endhead
\bottomrule
\endfoot
\bottomrule
\endlastfoot

Sjöman et al.\ \cite{SGJ23}
& Cross-sectional survey; $n=434$ Finnish climbers
& Finger (SPIIF onset)
& Association of regular finger training (RFT) with sport-impairing injury by experience level
& $\geq$6\,yr: $\chi^2=4.57$, $p=0.03$ (protective); $<$6\,yr, $\geq$7a: $\chi^2=4.61$, $p=0.03$ (risk)
& \textbf{F} (injury incidence); observational; no causal inference
& Moderate \\

\addlinespace

Klauser et al.\ \cite{KFB+02}
& Diagnostic accuracy; $n=64$ climbers, 75 fingers
& Pulley injuries (A2/A3/A4)
& Dynamic US sensitivity/specificity vs.\ MRI reference
& Sensitivity 98\%, specificity 100\%; verification bias acknowledged
& \textbf{F} (diagnostic accuracy)
& Moderate \\

\addlinespace

Berrigan et al.\ \cite{BWC+21}
& SR imaging review
& A2 pulley diagnosis
& US vs.\ MRI/CT sensitivity/specificity
& Sensitivity/specificity 0.94--1.00
& \textbf{F} (diagnostic accuracy)
& Moderate; secondary synthesis \\

\addlinespace

Sch\"{o}ffl et al.\ \cite{SHWS03}
& Prospective registry cohort; $n=604$ climbers, 122 pulley injuries
& Pulley injuries Grade I--IV
& Return to prior climbing level after conservative vs.\ surgical management
& All conservative Gr.~I--III eventually returned; 7 had prolonged recovery due to persistent tenosynovitis; Grade IV surgical outcomes largely good; no controls
& \textbf{F} (return to sport); no ES; no RCT
& Low--moderate \\

\addlinespace

Schneeberger \& Schweizer \cite{SS16}
& Prospective case series; $n=47$ ruptures in 45 climbers
& Complete pulley ruptures (Gr.~II--IV)
& TBD reduction and return to prior climbing level with PPS
& A2 TBD: 4.4$\pm$1.0\,mm $\to$ 2.3$\pm$0.6\,mm (50\% reduction); 38/43 (88.4\%) regained prior level; mean return 8.8\,months
& \textbf{F} (TBD surrogate + return to sport)
& Moderate (prospective case series; no controls) \\

\addlinespace

Mohn et al.\ \cite{MSM+22}
& Retrospective descriptive; $n=65$ climbers
& Flexor tenosynovitis
& Return to climbing and symptom duration
& All returned to climbing; ${\approx}$75\% regained or exceeded prior level; mean symptom duration 30.5\,wk; outcomes independent of specific therapy content
& \textbf{F} (return to sport)
& Low--moderate (retrospective; no controls) \\

\addlinespace

Lutter et al.\ \cite{LSS+18}
& Consecutive case series; $n=60$ (57 climbers)
& Lumbrical muscle tear
& Return to climbing by grade; imaging classification
& All recovered with functional therapy; Grade III healing significantly longer; return 6$+$\,wk (Gr.~I--II), 3--6\,months (Gr.~III)
& \textbf{F} (return to sport); imaging by US/MRI
& Moderate (largest dedicated series; no controls) \\

\addlinespace

Sch\"{o}ffl et al.\ \cite{SLL+25}
& Prospective case series; $n=69$ capsulitis cases
& PIP joint capsulitis/synovitis
& Return to prior climbing level at 12\,months
& 72\% returned to prior level; 28\% showed decreased performance (mainly from extended rest, not residual symptoms); staged algorithm effective
& \textbf{F} (return to sport)
& Moderate (prospective case series; no controls) \\

\addlinespace

Schweizer \cite{Sch00}
& Biomechanical in-vivo; $n=16$ fingers
& A2 pulley loading / taping
& Bowstringing distance and force absorption under taping variants
& Proximal A2: bowstringing $\downarrow$2.8\%, force absorbed ${\approx}11\%$; distal: $\downarrow$22\%, ${\approx}12\%$
& \textbf{S} (biomechanical proxy; no clinical outcome)
& Moderate (biomechanics) \\

\addlinespace

Alfredson et al.\ \cite{APJL98}
& Prospective cohort $+$ historical controls; $n=15+15$
& Chronic Achilles tendinosis (analogue)
& Return to running at 12\,wk: eccentric heel-drop vs.\ continued activity
& 100\% return (eccentric) vs.\ 0\% (control midpoint); no formal ES
& \textbf{F} (return to sport); transfer to pulleys bioplausible only
& High (Achilles); ESS 3 pulley transfer \\

\addlinespace

Cook \& Purdam \cite{CP08}
& Conceptual review / mechanistic model
& Tendinopathy (general)
& Reactive$\to$dysrepair$\to$degenerative model; load-management framework
& No ES; conceptual only
& \textbf{S} (mechanistic model)
& Moderate; well-cited \\

\addlinespace

Rio et al.\ \cite{RKP+15}
& RCT crossover; $n=6$/condition (patellar tendinopathy)
& Patellar tendinopathy (analogue)
& Immediate pain (VAS) and MVIC: isometric vs.\ isotonic
& Pain: 7.0$\pm$2.04 $\to$ 0.17$\pm$0.41 (isometric); MVIC $\uparrow$18.7$\pm$7.8\%; no CI reported
& F (VAS pain); S (MVIC); patellar only
& Moderate; $n=6$; single session \\

\addlinespace

Shaw et al.\ \cite{SLBR+17}
& RCT crossover; $n=8$
& Connective tissue (ligament model)
& P1NP AUC; engineered-ligament mechanics post-supplement
& P1NP ${\approx}2{\times}$ placebo ($p{<}0.05$); ligament stiffness $\uparrow$; no CI reported
& \textbf{S} (P1NP; in-vitro ligament)
& Moderate; $n=8$ \\

\addlinespace

Iruretagoiena et al.\ \cite{IUDORL+20}
& SR; 7 studies
& A2/A4 pulley diagnosis
& US tendon--bone distance cutoff thresholds
& No consensus; A2 rupture cutoffs 1.9--5.1\,mm across studies
& \textbf{F} (diagnostic accuracy); threshold heterogeneity limits utility
& Moderate diagnostic SR \\

\end{longtable}
}% end footnotesize

%---------------------------------------------------------------------
\subsection{Injury-Risk Stratification Framework}
\label{subsec:riskstrata}

Risk tiers used throughout this section. Incidence estimates derive from available survey and case series data; where climbing-specific data are absent, estimates are flagged as extrapolated from adjacent sports.

\renewcommand{\arraystretch}{1.3}
\begin{center}
\small
\begin{tabular}{>{\raggedright}p{2.4cm}
                >{\raggedright}p{3.0cm}
                >{\raggedright}p{4.0cm}
                >{\raggedright\arraybackslash}p{5.0cm}}
\toprule
\textbf{Stratum} & \textbf{Estimated incidence} & \textbf{Definition} & \textbf{Management implication} \\
\midrule
\textbf{Low} & $<$5\,\%/year & Managed with standard warm-up and load monitoring; no elevated structural risk & Routine load progression; standard 48\,h recovery; no additional screening \\
\addlinespace
\textbf{Moderate} & 5--20\,\%/year & Context-dependent risk; structural risk present but mitigable with specific precautions & Explicit monitoring of finger pain VAS; HRV supplementary tracking; grip-position modification if VAS $\geq$2/10 \\
\addlinespace
\textbf{High} & $>$20\,\%/year & Elevated structural risk; explicit mitigation required & Mandatory load reduction, grip-position restriction, structured rehabilitation; medical evaluation if VAS $\geq$3/10 at rest \\
\addlinespace
\textbf{Very High} & Near-certain injury without intervention & Contraindicated without medical supervision & Full climbing rest pending medical evaluation; return protocol requires clinical clearance \\
\bottomrule
\end{tabular}
\end{center}

\noindent\textbf{Note on incidence data.} No prospective climbing study has measured annual pulley injury incidence by stratum. The rate estimates above are informed by Sch\"{o}ffl et al.\ (2003, \cite{SHWS03}) and Sjöman et al.\ (2023) survey data and should be treated as indicative. Injury incidence and return-to-sport risk are separate quantities: a Low-incidence injury can carry High return-to-sport risk if mismanaged. These strata are stated explicitly at each recommendation below where applicable.

%---------------------------------------------------------------------
\subsection{A2 Pulley Injuries}
\label{subsec:a2pulley}

\subsubsection{Anatomy and Mechanism}

The flexor tendon sheath of digits II--V contains five annular (A1--A5) and three cruciate (C1--C3) pulleys. The A2 pulley spans the proximal half of the proximal phalanx and is the largest, least deformable annular pulley; it prevents bowstringing of the flexor digitorum profundus (FDP) and flexor digitorum superficialis (FDS) by maintaining the tendons' mechanical line of action. In published case series, the ring (digit IV) and middle (digit III) fingers account for $\approx$98\% of pulley injuries in climbers, consistent with their role as the primary force-bearing digits in crimp configurations.

\begin{bioplausbox}{Full-Crimp Loading Hierarchy and Eccentric Failure Mechanism · \bioplausible · ESS 2}
\textbf{Grip-position loading hierarchy (Bioplausible):} Full crimp (PIP $\approx$90$^\circ$ flexion, DIP hyperextension or 0--20$^\circ$) $>$ half crimp $>$ open hand, due to acute tendon approach angle against the pulley increasing contact pressure. Schweizer (2001) estimated full-crimp A2 forces at 35--40\,N/kg body mass in cadaveric/model studies, half-crimp at $\approx$30--35\,N/kg, and open-hand at $\approx$20--25\,N/kg. No 95\,\% CI was reported. The Vigouroux et al.\ (2006, \cite{VQLVM06}) biomechanical model reports A2 forces $\approx$36$\times$ higher in crimp versus slope grip, and A4 forces $\approx$4$\times$ higher --- a substantially larger differential than the Schweizer absolute-force figures alone convey.

Sch\"{o}ffl I.\ et al.\ (2009, \cite{SOJ+09}) cadaveric study of concentric vs.\ eccentric loading found \textit{eccentric} pulley stress (sudden finger extension under load, as in a foot slip or a hold crimped open by surprise) to be the dominant failure mechanism, not static contact pressure alone. These data support recommending open-hand grip during volume accumulation and caution during eccentric-loaded holds; however, no prospective injury-prevention RCT has confirmed that grip-position prescription reduces pulley injury incidence in climbers. \textbf{[Bioplausible: mechanistic from Schweizer 2000, Vigouroux 2006 \cite{VQLVM06}, and Sch\"{o}ffl 2009 \cite{SOJ+09}; not confirmed in prospective trial.]}
\end{bioplausbox}

\subsubsection{Grading}

The Sch\"{o}ffl (2003) grading system \cite{SHWS03} remains the dominant clinical and research classification:

\begin{itemize}[leftmargin=1.5em]
  \item \textbf{Grade I}: Pulley strain without significant tendon--bone dehiscence; TBD increase $<$1--2\,mm on dynamic US. Tenderness over pulley; pain with resisted PIP flexion.
  \item \textbf{Grade II}: Complete rupture of A4 pulley \textit{or} partial rupture of A2 or A3 pulley. Mild or absent bowstringing; palpable swelling; pain with active loading.
  \item \textbf{Grade III}: Complete rupture of A2 \textit{or} A3 pulley. Clear bowstringing (TBD $>$2--3\,mm under load on US); may have bruising and palpable defect.
  \item \textbf{Grade IVa}: Double-pulley rupture (A2/A3 or A3/A4). Substantial bowstringing; significant functional deficit (Simon et al.\ 2022, \cite{SLTS22}).
  \item \textbf{Grade IVb}: Triple-pulley rupture (A2/A3/A4) \textit{or} single/double rupture with complicating pathology: Flap Irritation Phenomenon (FLIP), progressive contracture, tenosynovitis, collateral ligament injury, or lumbrical involvement (\cite{SLTS22}).
\end{itemize}

\begin{bioplausbox}{Schöffl Grading System: Clinical Utility Without RCT Validation · \bioplausible · ESS 2}
\textbf{Grading validity (Bioplausible):} The Sch\"{o}ffl system derives from case series, not a validated prospective grading-outcome study. No inter-rater reliability data from RCTs exists. Grade IV was refined by Simon et al.\ (2022, \cite{SLTS22}) into IVa (double-pulley) and IVb (triple-pulley or with complications), which has direct surgical-threshold implications. The system is clinically useful for mapping structural severity to management pathway but should not be treated as a validated prognostic instrument. \textbf{[Bioplausible: supported by case series \cite{SHWS03}; Grade IV refinement from Simon et al.\ 2022 \cite{SLTS22}; no RCT validation.]}
\end{bioplausbox}

\subsubsection{Diagnosis}

\paragraph{Physical examination.}
Pain reproduced by resisted PIP flexion and direct palpation over the proximal volar aspect of the proximal phalanx is the primary clinical indicator. Bowstringing may be palpable or visible in Grade III--IV when the digit is actively flexed against resistance. Sensitivity and specificity of clinical examination alone are not established in controlled diagnostic accuracy studies.

\begin{bioplausbox}{Clinical Examination Reliability for Pulley Injury: Unvalidated · \bioplausible · ESS 2}
\textbf{Clinical exam reliability (Bioplausible):} Focal tenderness and pain with resisted PIP flexion are widely used but unvalidated against imaging or surgical reference standards in controlled studies. \textbf{[Bioplausible: clinical reasoning; no ROC/diagnostic accuracy study.]}
\end{bioplausbox}

\paragraph{Imaging.}\leavevmode

\begin{evidencebox}{Ultrasound as Primary Diagnostic Modality for Pulley Injury · \evidencebacked · ESS 6}
\textbf{Ultrasound as primary modality (Evidence-backed):} Dynamic high-resolution US ($\geq$15\,MHz linear probe, assessed in passive extension and full resisted flexion) is the recommended first-line modality. Klauser et al.\ (2002) reported US sensitivity 98\%, specificity 100\% vs.\ MRI reference standard in 64 climbers with 75 symptomatic fingers (\cite{KFB+02}; verification bias acknowledged). Berrigan et al.\ (2021) systematic review confirms US sensitivity/specificity 0.94--1.00 across included studies (\cite{BWC+21}). Dynamic assessment comparing TBD at 20$^\circ$ passive extension vs.\ full active flexion is essential; static views underestimate Grade II bowstringing and miss partial tears.

Complete A2 rupture TBD diagnostic cutoffs are inconsistent across studies, ranging 1.9--5.1\,mm (Iruretagoiena et al.\ 2020, \cite{IUDORL+20}). Miro et al.\ (2021, \cite{MVSS21}) found that experienced high-level climbers have physiologically elevated TBD (0.3--0.4\,mm longer at A2/A4) versus non-climbers \textit{without reported injury}, indicating that TBD thresholds derived from non-climbing populations may over-diagnose injury in adapted athletes. MRI sensitivity for complete pulley tears ranges 0.79--1.00 across studies in varying finger positions; MRI is superior for multi-structure injuries (capsulitis, collateral ligament, volar plate). \textbf{Bottomline: use high-resolution dynamic US first; adjust TBD thresholds for climber adaptation; reserve MRI for equivocal presentations or suspected multi-structure Grade IVa/IVb injury.}
\end{evidencebox}

%---------------------------------------------------------------------
\subsubsection{Conservative Management --- Grade I--II and Most Grade III}

\begin{warningbox}[title=No Climbing Rehabilitation RCTs]
All immobilisation durations, return-to-loading timelines, and hangboard progression thresholds below are derived from case series (Sch\"{o}ffl et al.\ 2003, \cite{SHWS03}) or extrapolated from non-climbing tendinopathy RCTs. None has been validated in a controlled climbing study.
\end{warningbox}

\paragraph{Immobilisation.}\leavevmode

\begin{bioplausbox}{Short-Term Immobilisation: Grade-Based Duration (Schöffl 2003) · \bioplausible · ESS 2}
\textbf{Short-term immobilisation (Bioplausible):} Typical recommendations from Sch\"{o}ffl et al.\ (2003, \cite{SHWS03}) case series:
\begin{itemize}[leftmargin=1.5em,noitemsep]
  \item Grade I: 0\,d immobilisation; functional therapy 2--4\,wk; tape protection.
  \item Grade II: 10\,d immobilisation with pulley-protection splint (PPS); then 2--4\,wk functional therapy.
  \item Grade III: 14\,d immobilisation with PPS; PPS worn continuously for 6--8\,wk total; functional therapy alongside.
\end{itemize}
Rationale: early controlled mechanical stimulation promotes aligned collagen deposition; prolonged immobilisation ($>$1\,wk) is inconsistent with tendon and ligament healing biology. However, no climbing RCT has compared immobilisation durations. \textbf{[Bioplausible: extrapolated from tendon biology; case series source.]}
\end{bioplausbox}

\begin{evidencebox}{Pulley-Protection Splint (PPS): Primary Conservative Device for Grade II--III · \evidencebacked · ESS 5}
\textbf{Pulley-protection splint (Evidence-backed for conservative outcomes):} The PPS (thermoplastic circumferential ring worn full-time) physically approximates the tendon to the bone continuously, reducing TBD passively at rest and during activity. Schneeberger \& Schweizer (2016, \cite{SS16}) prospective series ($n=47$ complete ruptures in 45 climbers):
\begin{itemize}[leftmargin=1.5em,noitemsep]
  \item PPS worn continuously for 2\,months; active functional therapy alongside from early weeks.
  \item A2 TBD reduced from 4.4$\pm$1.0 to 2.3$\pm$0.6\,mm (50\% reduction); A4 TBD from 2.9$\pm$0.7 to 2.1$\pm$0.5\,mm.
  \item 38/43 climbers (88.4\%) regained prior climbing level at mean 8.8\,months.
\end{itemize}
The Sch\"{o}ffl (2020, \cite{SSF+20}) MLTJ review states: ``The pulley support ring therapy must be started within a few days after trauma and performed strictly for 6--8 weeks.'' PPS and H-tape serve complementary roles: PPS provides continuous passive bowstringing reduction during healing; H-tape provides additional support during sport. \textbf{[Evidence-backed for TBD reduction and functional return: prospective case series \cite{SS16}; no RCT comparison vs.\ H-tape or immobilisation only.]}
\end{evidencebox}

\paragraph{Return-to-loading timeline.}\leavevmode

\begin{bioplausbox}{Return-to-Climbing Timelines After Pulley Injury · \bioplausible · ESS 2--5}
\textbf{Return-to-climbing timelines:} Sch\"{o}ffl et al.\ (2003, \cite{SHWS03}) case series provides the foundational reference:
\begin{itemize}[leftmargin=1.5em,noitemsep]
  \item Grade I: progressive reloading at 1--2\,wk; return to climbing 4--6\,wk.
  \item Grade II: progressive reloading at 3--4\,wk; return to climbing 6--10\,wk.
  \item Grade III (conservative with PPS): progressive reloading at 6--8\,wk; Schneeberger \& Schweizer (2016, \cite{SS16}) prospective PPS series reports mean 8.8\,months to prior climbing level --- substantially longer than the Sch\"{o}ffl (2003) ``3--5\,months'' figure, which derived from a subset without continuous PPS. Grade III patients should be counselled on 6--10\,months as a realistic expectation.
\end{itemize}
\textbf{[Grade I--II: Bioplausible/Folklore from Sch\"{o}ffl 2003 \cite{SHWS03}. Grade III: ESS 5 prospective case series \cite{SS16}; mean 8.8\,months.]}
\end{bioplausbox}

\paragraph{H-taping and ring splinting.}\leavevmode

\begin{bioplausbox}{H-Taping: Biomechanical Evidence and Limitations · \bioplausible · ESS 2}
\textbf{Biomechanical basis and limitations (Bioplausible):} Schweizer (2000, \cite{Sch00}) in-vivo study ($n=16$ fingers) found tape placed distally over the proximal phalanx (the more effective position) reduced bowstringing by 22\% and absorbed $\approx$12\% of bowstringing force; more proximal placement over the A2 pulley gave only 2.8\% reduction. Sch\"{o}ffl et al.\ (2007, \cite{SES+07}) controlled biomechanical study ($n=8$ climbers) found H-tape reduced TBD by 16\% and increased crimp MVC by $+$13\% (no CI) vs.\ no tape; no effect in open-hand hang.

A systematic review by Larsson et al.\ (2022, \cite{LNB22}) reported low-to-moderate certainty for bowstringing reduction (15--22\%) in \textit{injured} fingers only; no significant effect in uninjured fingers; zero RCTs on re-injury. Salas et al.\ (2022, \cite{SMT+23}) cadaveric biomechanical study (56 paired fresh-frozen specimens) found H-taping did \textit{not} significantly increase failure force or reduce bowstringing in partial A2 tears, and provided no protection against A2 rupture in intact specimens. The cadaveric findings directly challenge H-tape's biomechanical protective rationale; clinical utility may rest primarily on pain modulation, proprioception, and confidence. \textbf{[Bioplausible: positive small-sample data from \cite{Sch00,SES+07}; contradicted by cadaveric evidence \cite{SMT+23}; no RCT for clinical outcomes.]}
\end{bioplausbox}

\begin{folklorebox}{H-Taping in Climbing Rehabilitation: Universal Practice · \folkloreverified · ESS 1}
\textbf{H-taping clinical practice (Folklore):} H-taping around the proximal phalanx (figure-8 configuration creating dorsal bridge) is near-universal in climbing rehabilitation practice. Community consensus, Hooper's Beta, Schöffl-era clinical guidelines. No RCT or controlled case series has measured re-injury rates, healing time, or functional outcomes with vs.\ without taping. \textbf{[Folklore: climbing community / climbing physio consensus.]}
\end{folklorebox}

\paragraph{Load-managed rehabilitation --- tendinopathy extrapolation.}\leavevmode

\begin{evidencebox}{Tendon Loading Paradigms: Eccentric and Isometric Evidence (Achilles/Patellar) · \evidencebacked · ESS 3}
\textbf{Tendon loading paradigms (Evidence-backed for analogous tissues):}
\begin{itemize}[leftmargin=1.5em]
  \item Alfredson et al.\ (1998, \cite{APJL98}): heavy-load eccentric heel-drop (3$\times$15 twice daily, 12\,wk) achieved 100\% return to running in $n=15$ chronic Achilles tendinosis patients vs.\ 0\% for continued activity alone. A 2023 meta-analysis of 5 Achilles RCTs reported eccentric exercise pain reduction MD $-$1.21 (95\,\% CI $-$2.72 to $-$0.30) vs.\ other exercise.
  \item Cook \& Purdam (2009, \cite{CP08}): reactive--dysrepair--degenerative continuum model; reactive stage warrants load reduction (not rest) and isometric loading; degenerative stage benefits from progressive isotonic loading.
  \item Rio et al.\ (2015, \cite{RKP+15}): isometric leg extension produced immediate pain reduction of 6.8/10 (from 7.0$\pm$2.04 to 0.17$\pm$0.41, sustained 45\,min) and MVIC $\uparrow$18.7$\pm$7.8\% vs.\ isotonic in patellar tendinopathy ($n=6$).
\end{itemize}
\textbf{[Evidence-backed for Achilles and patellar tendinopathy. Extrapolation to A2 pulley is Bioplausible only --- see below.]}
\end{evidencebox}

\begin{bioplausbox}{Progressive Loading Applied to A2 Pulley: Extrapolation from Tendinopathy · \bioplausible · ESS 2}
\textbf{Pulley-specific loading (Bioplausible):} The A2 pulley is a fibrocartilaginous ligamentous condensation of the flexor sheath --- primarily type\,I collagen with fibrocartilaginous zones under compressive load --- not a tendon. Cellular overlap with tendons exists (fibroblast mechanosensing, integrin-mediated collagen remodelling), but load-response properties are not identical. Application of progressive loading protocols (isometrics $\to$ isotonic $\to$ functional) to pulley rehabilitation is bioplausible but constitutes an extrapolation from a different tissue type. No peer-reviewed study has measured A2 pulley load during hangboard exercises with validated force-sensing systems and correlated measurements with tissue healing outcomes. \textbf{[Bioplausible: mechanistic framework from Cook \& Purdam 2009 and Alfredson 1998; Folklore: specific thresholds and protocols.]}
\end{bioplausbox}

\paragraph{Load-managed hangboard rehabilitation during recovery.}\leavevmode

\begin{folklorebox}{Tyler Nelson Progressive Loading Protocols for Pulley Rehabilitation · \folkloreverified · ESS 1}
\textbf{Tyler Nelson / Camp 4 Human Performance protocols (Folklore):} Tyler Nelson has disseminated progressive loading protocols for pulley rehabilitation through clinical practice and online/professional channels. Proposed structure: (a) isometric hangs beginning at 20--40\% perceived maximum in the early reactive phase, pain $\leq$2--3/10; (b) progress load 5--10\% per week guided by symptom response; (c) grip-position progression --- open-hand $\to$ half crimp, full crimp avoided until $\approx$week 12. These thresholds are not empirically derived from pulley-specific loading studies. \textbf{[Folklore: Tyler Nelson, C4HP, practitioner/coaching channels. Mechanistic framework Bioplausible via Cook \& Purdam 2009; specific thresholds unvalidated.]}
\end{folklorebox}

%---------------------------------------------------------------------
\subsubsection{Surgical Management --- Grade III--IV}

\begin{bioplausbox}{Surgical Indications for Pulley Injury · \bioplausible · ESS 2--3}
\textbf{Surgical indications (Bioplausible / case series):} Current consensus (Simon et al.\ 2022, \cite{SLTS22}; Schneeberger \& Schweizer 2016, \cite{SS16}) classifies Grade I--III as conservative first-line. Surgery is indicated for:
\begin{itemize}[leftmargin=1.5em,noitemsep]
  \item Grade IVa or IVb with persistent bowstringing or functional deficit after adequate conservative management.
  \item Flap Irritation Phenomenon (FLIP): pulley stump folding under the intact sheath, producing constant mechanical friction; US shows oscillating hyperechoic structure; MRI confirms. Resection $\pm$ repair resolves (Sch\"{o}ffl et al.\ 2011, \cite{SBS11}).
  \item Grade III with documented failure of $\geq$6--9\,months of compliant PPS + functional therapy.
\end{itemize}
Techniques: graft-based loop reconstruction (palmaris longus, extensor retinaculum) remains the standard for Grade IVa/b multi-pulley loss. Direct suture repair without graft (Balandier et al.\ 2026, \cite{BBM+26}) was reported in 8 elite climbers (6 isolated, 2 multiple ruptures): mean return to pre-injury level 5.8\,months (range 3--8), US-confirmed pulley continuity at 3\,years --- a substantially shorter timeline than graft reconstruction but limited to small-$n$ series.

Arora et al.\ (2007) retrospective series ($n=23$) reported excellent/good Buck-Gramcko outcomes in 21/23.\singlesource Bouyer et al.\ (2016) retrospective series ($n=38$, mean FU 85\,months)\singlesource found post-operative US E-space $<$2\,mm predicted return in 70\% vs.\ 25\%. Return-to-climbing: 3--6\,months (Grade III failed conservative, direct suture); 6--12\,months (Grade IVa/b graft).

\textbf{[Bioplausible / Low-quality evidence: all data from case series; no controls; heterogeneous techniques. Grade III should not receive first-line surgery \textemdash{} Schneeberger \& Schweizer \cite{SS16} showed 88.4\% return with conservative PPS even for complete ruptures.]}
\end{bioplausbox}

%---------------------------------------------------------------------
\subsection{Flexor Tendon Strains}
\label{subsec:flexortendon}

\begin{bioplausbox}{Flexor Tendon Strain: Anatomy and Mechanism · \bioplausible · ESS 2}
\textbf{Anatomy and mechanism (Bioplausible):} FDP inserts distally to the distal phalanx; FDS splits at the proximal phalanx to insert bilaterally onto the middle phalanx base. Both run through the same fibrous sheath as the A2 pulley. Flexor tendon strains (partial or microtear injuries to the tendon proper) typically occur at the musculotendinous junction or within the tendon substance under repeated high-load crimp grip loading or sudden eccentric overload. They co-occur with Grade IV pulley injuries (definitional) or may present independently. No climbing-specific epidemiological data distinguish flexor tendon strain prevalence from pulley injury prevalence in isolation. \textbf{[Bioplausible: anatomical reasoning and clinical report.]}
\end{bioplausbox}

\paragraph{Grading and differentiation from pulley injury.}\leavevmode

\begin{bioplausbox}{Flexor Tendon Strain Grading and Differentiation from Pulley Injury · \bioplausible · ESS 2}
\textbf{Grading (Bioplausible / clinical consensus):}
\begin{itemize}[leftmargin=1.5em,noitemsep]
  \item Grade I: minor fibre disruption; no architectural loss on US; tenderness along tendon course; pain with resisted flexion.
  \item Grade II: partial tear; hypoechoic defect on US; weakened active flexion; focal swelling.
  \item Grade III: complete/near-complete disruption; functional deficit; requires urgent surgical referral.
\end{itemize}
Distinguishing from pulley injury: pulley injury presents with focal tenderness over the proximal phalanx with bowstringing on dynamic US; flexor tendon strain presents with pain along the tendon course proximal or distal to the pulley, internal tendon changes on US, but no bowstringing. Precise attribution requires high-quality dynamic US. \textbf{[Bioplausible: no climbing-specific grading validation study.]}
\end{bioplausbox}

\paragraph{Management.}\leavevmode

\begin{bioplausbox}{Grade I--II Flexor Tendon Strain: Conservative Management · \bioplausible · ESS 2}
\textbf{Grade I--II management (Bioplausible):} Principles mirror A2 pulley Grade I--II management: (a) acute load reduction 5--14\,d; (b) early progressive loading from week 2--3, pain-guided; (c) progressive return to full loading over 6--12\,wk. Mechanistic basis is identical to the Cook--Purdam continuum framework. No climbing-specific RCT or case series specifically distinguishes flexor tendon strain outcomes from pulley injury outcomes. \textbf{[Bioplausible: same extrapolation from tendinopathy literature.]}
\end{bioplausbox}

Grade III FDP avulsion (``jersey finger'' equivalent): requires urgent surgical referral and formal hand surgery protocol --- outside scope of conservative climbing rehabilitation.

%---------------------------------------------------------------------
\subsection{Flexor Tenosynovitis}
\label{subsec:tenosynovitis}

Flexor tenosynovitis --- inflammation of the tendon sheath surrounding the FDP/FDS tendons --- is the second most common climbing finger injury (7\% of all injuries in Sch\"{o}ffl 2003 \cite{SHWS03} and Lutter 2020 \cite{SSF+20} cohorts). It is the injury most frequently confused with pulley injury; mismanagement (treating as a pulley strain) prolongs recovery.

\begin{evidencebox}{Flexor Tenosynovitis: Diagnosis and Conservative Outcomes · \evidencebacked · ESS 5}
\textbf{Diagnosis:} US shows the \textit{halo phenomenon} --- peritendinous fluid in the tendon sheath around the flexor tendon. No bowstringing. Pain is diffuse along the tendon course (proximal to the pulley or radiating into the palm) rather than the focal proximal-phalanx tenderness of a pulley injury. Active glide against resistance reproduces pain.

\textbf{Natural history and conservative management} (Mohn et al.\ 2022, \cite{MSM+22}; retrospective $n=65$ climbers with confirmed tenosynovitis):
\begin{itemize}[leftmargin=1.5em,noitemsep]
  \item All patients returned to climbing; ${\approx}$75\% regained or exceeded prior climbing level.
  \item Mean symptom duration 30.5\,wk; course was undulating (flare-ups common).
  \item Load reduction was the dominant treatment (91\% of patients); no specific therapy content significantly predicted outcome in regression analysis, suggesting a favourable natural history.
  \item Predictors of longer recovery: high pain during daily activities at presentation; prior climbing grade $\geq$7b.
\end{itemize}

\textbf{Corticosteroid injection for chronic tenosynovitis} ($>$6\,wk; Sch\"{o}ffl et al.\ 2019, \cite{SSL19}; uncontrolled case series $n=42$): two injections 7--10\,days apart; 73.8\% pain-free after second injection; mean 84.2\% pain reduction; no complications. No surgery required in any case.

\textbf{[Evidence-backed for natural history and corticosteroid utility: case series \cite{MSM+22,SSL19}; no RCT; no control group.]}
\end{evidencebox}

%---------------------------------------------------------------------
\subsection{PIP Joint Capsulitis and Synovitis}
\label{subsec:capsulitis}

PIP joint capsulitis/synovitis is the third most common climbing finger injury (6--7\% of all injuries; $\approx$18.8\% of finger injuries per Sch\"{o}ffl 2020 \cite{SSF+20}). It is frequently misdiagnosed as a pulley injury because both present with PIP-region pain.

\begin{evidencebox}{PIP Capsulitis/Synovitis: Mechanism, Diagnosis, and Treatment Algorithm · \evidencebacked · ESS 5}
\textbf{Mechanism:} Repeated crimp loading generates high peak pressure in the PIP joint, driving inflammatory mediators, synovial effusion, and progressive capsular fibrosis. Untreated chronic capsulitis leads to early osteoarthritis; long-term climbers commonly show PIP osteophytes on imaging.

\textbf{Differential from pulley injury:} Diffuse PIP joint swelling and morning stiffness without bowstringing; pain on active flexion/extension throughout the arc (not just at load); absent focal tenderness over the proximal phalanx. US shows intra-articular effusion $\pm$ periarticular hyperemia.

\textbf{Staged treatment algorithm} (Sch\"{o}ffl et al.\ 2025, \cite{SLL+25}; prospective case series $n=69$; Tanguay et al.\ 2023, \cite{Vag23}):
\begin{itemize}[leftmargin=1.5em,noitemsep]
  \item \textbf{Stage 1}: Load reduction, anti-inflammatory (topical/oral short-course), contrast baths, joint mobilisation; avoid full crimp until resolved.
  \item \textbf{Stage 2} (failure of Stage 1 at 4--6\,wk): Corticosteroid or hyaluronic acid injection.
  \item \textbf{Stage 3} (chronic, refractory): Radiosynoviorthesis.
\end{itemize}
Outcome: 72\% returned to prior climbing level; 28\% showed decreased performance, predominantly attributed to prolonged forced rest rather than residual symptoms. All climbers resumed some climbing within 12\,months.

\textbf{[Evidence-backed: prospective case series \cite{SLL+25}; no RCT; no control condition.]}
\end{evidencebox}

%---------------------------------------------------------------------
\subsection{Lumbrical Muscle Injuries}
\label{subsec:lumbrical}

Lumbrical muscle tears are highly climbing-specific and are absent from most general sports medicine references. They are caused by the \textit{quadriga effect}: the 3rd and 4th lumbrical muscles have a unique bipennate dual-FDP origin; when the ring or middle finger is loaded in a pocket (flexed) while adjacent fingers are extended, a shearing force tears the lumbrical at its FDP attachment. Clinical presentation: pain in the \textit{palm} at or proximal to the MCP joint, radial side of the ring finger --- not over the proximal phalanx as in pulley injury.

\begin{evidencebox}{Lumbrical Tear: Diagnosis, Classification, and Outcomes · \evidencebacked · ESS 5}
\textbf{Diagnostic test (Lumbrical Stress Test):} Pain provoked by passive extension of the ring finger while the patient attempts to close their fist; pain disappears when the middle finger is simultaneously extended (neutralising the quadriga effect). This test is specific to lumbrical involvement.

\textbf{Classification} (Lutter et al.\ 2018, \cite{LSS+18}; $n=60$ consecutive patients, 57 climbers):
\begin{itemize}[leftmargin=1.5em,noitemsep]
  \item Grade I: microtrauma, US-negative; Grade II: muscle-fibre disruption visible on US; Grade III: musculotendinous disruption (US or MRI required).
\end{itemize}
\textbf{Outcomes:} All 60 patients recovered with adapted functional therapy. Grade III healing significantly longer than Grade I/II. Return to climbing: $\geq$6\,wk (Gr.~I--II); 3--6\,months (Gr.~III).

\textbf{Management:} Avoid single-finger pocket holds and differential-loading configurations during recovery. Progressive open-hand loading on matched holds; gradual reintroduction of pockets guided by pain response. Schweizer (2003, \cite{Sch03}) first described the clinical entity in three case reports.

\textbf{[Evidence-backed for natural history and classification: largest dedicated series \cite{LSS+18} ($n=60$); no RCT; no control.]}
\end{evidencebox}

%---------------------------------------------------------------------
\subsection{Adolescent Physeal Injuries}
\label{subsec:adolescent}

\begin{warningbox}{Adolescent Climbers: Pulley Framework Does Not Apply}
In climbers under $\approx$16--17 years (male) or $\approx$14--15 years (female), the PIP joint growth plate physis is the mechanical weak point, not the pulley. Crimp-grip loading preferentially stresses the open physis, producing \textit{primary periphyseal stress injuries} (PPSI). Presentation mimics a pulley strain (PIP swelling, pain with loading) but X-ray $\pm$ MRI is required. Sch\"{o}ffl et al.\ (2022, \cite{SSF+21}) prospective multicenter algorithm study ($n=$ multiple centres): stop climbing immediately; X-ray (may appear normal); MRI if X-ray negative and high suspicion; 6--12\,wk rest from climbing. Sch\"{o}ffl et al.\ (2025 BMJ Open) found physeal injuries in 29.9\% of 137 injuries across 95 adolescent patients. Applying adult pulley rehabilitation protocols to adolescents risks permanent growth-plate damage and joint deformity.
\end{warningbox}

%---------------------------------------------------------------------
\subsection{Skin Management}
\label{subsec:skin}

\subsubsection{Flappers --- Acute Traumatic Skin Tears}

A flapper is a traumatic partial-thickness skin tear, occurring at the callus-to-normal-skin transition where a stress riser under tangential shear load on rock hold edges causes delamination. Most common on the second and third digit palmar skin and thenar eminence.

\begin{folklorebox}{Acute Flapper Management: Universal Climbing Practice · \folkloreverified · ESS 1}
\textbf{Acute management (Folklore):} Trim any non-viable detached skin flap with sterile scissors (reduces bacterial trapping and maceration); irrigate; apply occlusive dressing or thin skin balm (petrolatum, Joshua Tree Salve, Climb On, or similar humectant) to maintain moist wound environment. Taping over a partially healed wound allows earlier return. \textbf{[Folklore: climbing community universal practice. General wound healing biology supports moist wound environment for epidermal regeneration (bioplausible extrapolation), but no climbing-specific trial exists.]}
\end{folklorebox}

\begin{toverifybox}{Return-to-Climbing Timeline After Flappers · \toverify · ESS 1}
\textbf{\toverify:} Anecdotal community consensus: mild flapper (superficial, $<$1\,cm$^2$) 3--7\,d; moderate (deeper or larger) 7--14\,d; larger raw dermis 1--2\,wk --- dependent on pain, infection risk, and hold texture. No empirical healing timeline data exist. Requires $\geq$3 independent named climbing sources with specific timelines for \folkloreverified status.
\end{toverifybox}

\begin{folklorebox}{Callus Management for Flapper Prevention · \folkloreverified · ESS 1}
\textbf{Prevention --- callus management (Folklore):} Filing calluses flush (reducing height and smoothing the callus-to-normal-skin transition) is near-universal in elite climbing practice; mechanistic goal is to reduce the stress riser at the thickness transition. Has not been evaluated in any study measuring flapper incidence as a function of callus management. Additional community recommendations: moisturise nightly, avoid alcohol-based sanitisers before climbing, manage session volume on sharp holds, rotate grip styles. \textbf{[Folklore: Eric Horst, \textit{Training for Climbing}; Lattice Training / Will Anglin; athlete and coach consensus.]}
\end{folklorebox}

\subsubsection{Chronic Skin Thinning}

\begin{bioplausbox}{Chronic Skin Thinning: Mechanism from Repeated Hold Friction · \bioplausible · ESS 2}
\textbf{Mechanism (Bioplausible):} Repeated friction on small-radius holds reduces stratum corneum thickness below functional baseline, presenting as hypersensitivity, split skin, and reduced friction coefficient. Contributing factors: training volume, hold substrate roughness, ambient humidity (very low humidity accelerates transepidermal water loss and skin cracking), and individual skin phenotype. No climbing-specific study has measured stratum corneum thickness longitudinally as a function of training load. \textbf{[Bioplausible: mechanism from general dermatology; no climbing-specific intervention data.]}
\end{bioplausbox}

\begin{folklorebox}{Climbing Skin Products and Humidity Management · \folkloreverified · ESS 1}
\textbf{Interventions (Folklore):} Climbing-specific skin products (Rhino Skin Solutions, ClimbSkin, Friction Labs products), humidity management (40--60\% indoor humidity anecdotally preferred), Antihydral (methenamine; acts by protein denaturation to reduce sweating, potentially increasing friction but risking ``glassy'' skin and splits), callus sanding, and session spacing are community-recommended. No peer-reviewed controlled trial has evaluated any climbing skin product for skin health, injury rate, or return-to-climbing outcomes. \textbf{[Folklore: manufacturer claims, community/practitioner consensus; Dave MacLeod, ClimbSkin. Emollient/humectant mechanisms are Bioplausible from general dermatology literature on skin barrier repair.]}
\end{folklorebox}

\subsubsection{Collagen Supplementation for Skin Healing}

\begin{evidencebox}{Collagen Supplementation: Skin and Connective Tissue Endpoints in Non-Climbing Populations · \evidencebacked · ESS 3}
\textbf{Collagen synthesis markers and skin endpoints (Evidence-backed for non-climbing populations):}
\begin{itemize}[leftmargin=1.5em]
  \item Shaw et al.\ (2017, \cite{SLBR+17}): RCT crossover ($n=8$); 15\,g gelatin + 48\,mg vitamin\,C, 1\,hr pre-exercise vs.\ placebo doubled circulating P1NP (collagen synthesis marker); in-vitro engineered ligament mechanics improved. \textit{Note: $n=8$ is very small; endpoints are surrogate; no injury or skin-healing outcome data.}
  \item A meta-analysis of oral hydrolyzed collagen peptides (26 RCTs, $n=1721$) for skin aging found improved hydration and elasticity endpoints. A systematic review/meta-analysis of 19 RCTs ($n=1341$) on oral/topical peptides found improved hydration and modest wrinkle reduction.
  \item Proksch et al.\ (2025, \cite{PZO25}): RCT ($n=66$), 2.5\,g collagen peptides vs.\ placebo --- reduced wrinkles, improved elasticity and hydration.
  \item A double-blind pilot RCT in burn patients ($n=30$): hydrolyzed collagen supplementation improved wound healing markers and reduced hospital stay $\approx$30\% vs.\ control.\singlesource
\end{itemize}
\textbf{[Evidence-backed for skin elasticity/hydration endpoints in aging and burn populations. Bioplausible for climbing skin healing: no direct evidence for flapper healing or friction-induced thinning in climbers.]}
\end{evidencebox}

%---------------------------------------------------------------------
\subsection{General Injury Prevention Principles}
\label{subsec:prevention}

\subsubsection{Warm-Up}

\begin{bioplausbox}{Warm-Up Efficacy for Injury Prevention: Mechanism Without Climbing RCT · \bioplausible · ESS 2}
\textbf{Warm-up efficacy (Bioplausible):} No climbing-specific RCT exists for injury prevention via warm-up. General sports medicine literature supports warm-up via increased connective tissue temperature (improving viscoelastic properties), increased neuromuscular activation, and enhanced motor pattern recruitment. The most methodologically robust proximate-domain evidence is the FIFA 11+ protocol (Soligard et al.\ 2008), which demonstrated $\approx$50\% lower-limb injury reduction in a large cluster-RCT in soccer --- tissue type, injury mechanism, and loading characteristics differ substantially from climbing finger pulley injuries.

Practical climbing warm-up: 5--10\,min aerobic movement; easy traversing; progressive finger recruitment; submaximal hangs before limit bouldering. \textbf{[Bioplausible: mechanistic basis exists; climbing-specific RCT evidence absent.]}
\end{bioplausbox}

\subsubsection{NSAIDs in Finger Injury Management}

\begin{bioplausbox}{NSAIDs for Finger Injuries: Contested Evidence, Pathology-Dependent Guidance · \bioplausible · ESS 2}
\textbf{NSAIDs (Bioplausible --- context-dependent):} NSAIDs are widely self-administered by climbers but their role in finger-injury healing is contested and pathology-dependent:
\begin{itemize}[leftmargin=1.5em,noitemsep]
  \item \textbf{Tenosynovitis / capsulitis (primarily inflammatory):} Short-course NSAIDs (48--72\,h acute phase) are consistent with anti-inflammatory management and broadly accepted in sports medicine. Evidence for clinical benefit in these conditions comes from general tendinopathy and synovitis literature, not climbing-specific trials.
  \item \textbf{Pulley / fibrocartilaginous tissue healing:} The evidence is conflicted. Heinemeier et al.\ (2017, \cite{HHM+17}) human tendinopathic tissue study found NSAIDs suppress COX-mediated prostaglandin E2 production that is required for tendon remodelling; the clinical significance at short doses is uncertain. Multiple in-vitro and animal models show COX-2 inhibition impairs tenocyte proliferation and collagen synthesis. However, no study has directly measured NSAID effects on A2 pulley healing in climbers.
  \item \textbf{Practical guidance:} Acetaminophen/paracetamol is preferred for analgesia during the first 2\,wk of acute pulley injury when bone-tendon healing is most active. If NSAIDs are used, limit to the initial 48--72\,h inflammatory phase only; avoid prolonged use during active tendon remodelling (weeks 1--6).
\end{itemize}
\textbf{[Bioplausible: mechanistic rationale from \cite{HHM+17}; no climbing-specific RCT on NSAID effects on pulley healing. Climbing community widely uses ibuprofen chronically through pain --- no outcome data exist.]}
\end{bioplausbox}

\subsubsection{Load Management --- The 10\% Rule}

\begin{toverifybox}{The 10\,\% Load Progression Rule Applied to Climbing · \toverify · ESS 1}
\textbf{\toverify --- 10\% rule (unvalidated in climbing):} The heuristic that weekly training load increases should not exceed 10\% of the prior week's load is widely cited in running medicine and climbing coaching. Its evidence base is weaker than commonly assumed: it derives primarily from retrospective epidemiological observation, not prospective RCT. Systematic reviews (e.g., Impellizzeri et al.\ 2020, \cite{IWC+20}) note that the specific 10\% threshold is arbitrary and inconsistently supported across sports. Gabbett (2016, \cite{Gab16}) and the acute:chronic workload ratio literature provide a more quantitative framework but remain primarily observational, contested in threshold values, and derived from team sports. For climbing, the threshold is additionally complicated because intensity, grip type, board angle, session density, and concurrent finger-training volume interact in ways not captured by a single linear loading metric. Requires $\geq$3 independent climbing-community sources citing this specific threshold for \folkloreverified status.
\end{toverifybox}

\subsubsection{Experience-Dependent Injury Risk}

\begin{evidencebox}{Experience-Dependent Finger Training Risk (Sjöman 2023) · \evidencebacked · ESS 5}
\textbf{Sjöman et al.\ (2023) --- Experience-dependent association (Evidence-backed):} Cross-sectional survey of $n=434$ Finnish climbers (\cite{SGJ23}) examined the association between route/bouldering-specific finger strength training (RFT) and sport/performance-impairing finger injury (SPIIF):
\begin{itemize}[leftmargin=1.5em,noitemsep]
  \item In climbers with $\geq$6\,yr experience: RFT \textbf{negatively} associated with SPIIF ($\chi^2=4.57$, $p=0.03$) --- structured training may be protective.
  \item In climbers with $<$6\,yr experience who had reached $\geq$7a (French sport): RFT \textbf{positively} associated with SPIIF ($\chi^2=4.61$, $p=0.03$) --- training in rapidly progressing novices may increase injury risk.
\end{itemize}
\textit{Quality constraints: cross-sectional design prohibits causal inference; self-reported injury and training data; Finnish climbing population; no dose-response data; no adjustment for confounders (training volume, hold type, warm-up, concurrent activities).}
\end{evidencebox}

\begin{bioplausbox}{Bidirectional Injury Risk Pattern: Mechanistic Interpretation · \bioplausible · ESS 2}
\textbf{Mechanistic interpretation (Bioplausible):} The bidirectional pattern is consistent with theoretical models of training-load injury risk in developing athletes (rapid capacity acquisition outpacing structural tissue adaptation) and the protective adaptation hypothesis in experienced athletes (accumulated structural adaptation of the fibrous flexor sheath, pulley hypertrophy, and neuromuscular co-contraction patterns). The mechanism for protection in experienced climbers has not been directly measured. \textbf{[Bioplausible for mechanism; Evidence-backed for the observed association in this dataset only.]}
\end{bioplausbox}

%---------------------------------------------------------------------
\subsection{What the Literature Cannot Yet Answer}
\label{subsec:injury-gaps}

\begin{enumerate}[leftmargin=2em]
  \item \textbf{Progressive loading vs.\ rest-only for Grade I--II A2 injuries.} Does structured hangboard rehabilitation outperform rest for time-to-return, re-injury rate, or ultrasound-confirmed healing in Grade I--II tears? No RCT exists; a moderate-sized trial ($n \geq 60$, Grade I--II, ultrasound-confirmed outcomes) is feasible.

  \item \textbf{Optimal immobilisation duration and type.} For Grade II--III A2 injuries, does 0--7\,d vs.\ 10--14\,d immobilisation (with standardised pulley protection) alter long-term TBD, return timeline, or re-injury rate? No controlled comparison exists.

  \item \textbf{H-taping clinical effectiveness.} Schweizer (2000, \cite{Sch00}) established small biomechanical reductions; Sch\"{o}ffl et al.\ (2007, \cite{SES+07}) reported acute crimp-strength improvement; Salas et al.\ (2022, \cite{SMT+23}) cadaveric study found no protective effect. Neither confirms reduced re-injury rate or faster healing in a controlled outcome study. Does H-taping change clinical outcomes, or does it act primarily via pain modulation and confidence? An RCT comparing PPS + H-tape versus PPS alone for clinical outcomes has not been conducted.

  \item \textbf{Ultrasound-guided return-to-climb criteria.} Does using serial TBD measurement to gate progression (e.g., TBD $<$2\,mm under resisted flexion before advancing grip position) reduce re-injury relative to symptom-only progression? The current TBD cutoff range of 1.9--5.1\,mm for complete rupture precludes standardised criteria.

  \item \textbf{Load-to-failure properties and healing kinetics of the human A2 pulley under progressive loading.} No in-vivo or ex-vivo study has quantified how incremental hangboard loading affects pulley tissue mechanics during healing. Tyler Nelson--style protocols lack any quantitative validation in the target tissue.

  \item \textbf{Collagen supplementation for pulley and skin healing.} Shaw et al.\ (2017) demonstrated elevated P1NP surrogate markers ($n=8$). No study has measured actual A2 pulley recovery time, re-injury rate, or imaging-based healing in climbers using gelatin + vitamin\,C. Similarly, no RCT has evaluated oral collagen peptides for flapper healing or chronic skin thinning under climbing-specific friction trauma.

  \item \textbf{Structural mediators of experience-dependent injury risk.} The Sjöman et al.\ (2023) bidirectional association is observational. A prospective cohort tracking novices from onset of structured finger training through 2--5\,yr, with serial ultrasound pulley morphology assessment, grip-force measurement, and load monitoring, would provide the first mechanistic evidence for whether pulley hypertrophy or neuromuscular adaptation mediates the protective effect observed in experienced climbers.

  \item \textbf{Pulley-protection splint vs.\ H-taping for clinical outcomes.} Schneeberger \& Schweizer (2016, \cite{SS16}) established PPS efficacy for TBD reduction and return-to-climbing, but no RCT has compared PPS + H-tape versus PPS alone versus H-tape alone for re-injury rate, return timeline, or patient-reported outcomes in Grade II--III injuries.

  \item \textbf{NSAID effects on pulley healing.} Climbers routinely self-medicate with ibuprofen through pulley pain. No study has measured the effect on healing outcomes in the A2 pulley (a fibrocartilaginous structure distinct from both tendon and bone). Whether short-course NSAIDs aid or harm healing in this specific tissue is unknown.

  \item \textbf{Tenosynovitis management: load reduction content vs.\ timing.} Mohn et al.\ (2022, \cite{MSM+22}) found no significant difference in outcomes based on specific therapy content; a prospective RCT comparing structured progressive loading versus unstructured load avoidance would clarify whether any specific rehabilitation element beyond ``reduce load'' adds clinical value.

  \item \textbf{Capsulitis natural history and OA progression.} The Sch\"{o}ffl (2025, \cite{SLL+25}) staged algorithm shows 72\% return to prior level, but long-term OA progression in those 28\% with decreased performance is not quantified. A 5--10\,yr cohort linking PIP synovitis episodes to imaging-confirmed OA development would establish whether aggressive early treatment is warranted.
\end{enumerate}

\input{sections/recovery}

%---------------------------------------------------------------------
\section{Conclusion}

\subsection*{Hangboarding}

Hangboarding is an evidence-supported supplemental training method for intermediate-to-advanced climbers. The meta-analytic evidence \cite{SO23} confirms consistent short-term improvements in dynamometric finger strength, RFD, and forearm endurance (SDM 0.41--1.23), but no trial has yet demonstrated grade improvement attributable to hangboarding alone. Max-hang protocols ($\geq$80\% MFS) best develop MFS; repeater protocols (60--80\% MFS) best develop stamina; 80\% MFS offers both. Daily low-intensity hangs (Abrahangs) produced a positive but non-significant trend versus climbing-only controls ($p = 0.12$) in the sole observational cohort; RCT confirmation is required before broad recommendation. Injury risk is experience-dependent: experienced climbers may benefit from a protective effect, while rapidly advancing less-experienced climbers face elevated incidence. Critical evidence gaps remain: female-cohort data, trials $>$10 weeks with injury tracking, and any functional on-wall grade outcome.

\subsection*{Weekly Training Integration}

Session sequencing and weekly structure are predominantly BIOPLAUSIBLE or FOLKLORE\_VERIFIED (ESS 1--2); no RCT has tested a complete weekly training template as a unit. The $\geq$48 h inter-session minimum between high-intensity finger sessions ($\geq$80\% MFS) is a study-design convention, not an empirically optimised prescription. Hangboard max-strength work must precede, not follow, limit bouldering when MFS is the target. Evidence for concurrent-training interference is well-established in lower-body strength-endurance models \cite{WMR+12} and is bioplausibly transferable to climbing, but no RCT has measured its magnitude in climbers. Deload indicators (RPE elevation, rolling force decrement, HRV trends) are extrapolated from general resistance training and have not been validated in climbing populations.

\subsection*{Supplementation}

No supplement in this review has been tested with a functional climbing outcome as the primary endpoint. Within that constraint: \textbf{collagen + vitamin C} (15\,g hydrolysed collagen $+$ 48\,mg ascorbic acid 1\,h pre-exercise) produces the most direct mechanism-to-tissue link for pulley and tendon support, though evidence remains surrogate (P1NP biomarker; ESS 3--4); \textbf{beta-alanine} (3.2--6.4\,g/day $\geq$4 weeks) has the strongest climbing-specific RCT signal for endurance tasks of 60--240\,s \cite{SNWK21} (ESS 7, but $n=15$); \textbf{caffeine} (3\,mg/kg) showed null results on all proxy climbing metrics in the only published climbing RCT \cite{RLMAMR+26} (ESS 7, $n=13$); \textbf{creatine} is evidence-backed for strength and power in general populations but carries a body-mass penalty that directly offsets performance on gravity-dependent movements; \textbf{omega-3} and \textbf{magnesium} lack climbing-specific controlled evidence and remain bioplausible at most for sub-populations with documented deficiency.

\subsection*{Injury Management}

A2 pulley injuries are the most common acute climbing finger injury; the corrected Sch\"{o}ffl Grade II--IV classification, pulley-protection splint (PPS) as primary conservative device (88.4\% return at mean 8.8\,months; \cite{SS16}), and Grade IVa/IVb surgical threshold are key updates from the initial framework. Flexor tenosynovitis (second most common; favourable natural history under load reduction, \cite{MSM+22}) and PIP capsulitis/synovitis (third most common; staged algorithm with 72\% return, \cite{SLL+25}) represent previously absent coverage that substantially broadens the section's clinical utility. Lumbrical tears are identified by the pathognomonic lumbrical stress test and managed with pocket-avoidance and progressive loading (\cite{LSS+18}). Adolescent physeal injuries (the dominant injury in $<$16-year-old climbers) require immediate rest and imaging; the adult pulley framework is contraindicated (\cite{SSF+21}). Injury risk is experience-dependent: a bidirectional pattern (elevated risk in intermediate climbers, apparent protection in highly experienced climbers) is observationally supported \cite{SGJ23} (ESS 5). Critical gaps include absence of RCTs comparing PPS versus H-taping, progressive loading versus rest-only for Grade I--II, NSAID effects on pulley healing, and long-term OA progression after capsulitis episodes.

\subsection*{Recovery}

No RCT has tested any recovery modality specifically against climbing performance or pulley/tendon injury incidence. Within that constraint: \textbf{sleep} $\geq$8\,h is the recovery modality with the strongest evidence base across sports \cite{MSB+14}; \textbf{cold water immersion} is evidence-backed for DOMS reduction but should be avoided immediately post-strength/hangboard sessions as it attenuates anabolic signalling \cite{RRM+15}; \textbf{massage and foam rolling} reduce perceived soreness (SMD $-$0.72; \cite{DDT+18}) but do not address tendon recovery on the collagen-remodelling timeline; \textbf{active recovery} (20--30\,min aerobic at RPE 3--4) accelerates lactate clearance and has one climbing-specific controlled study (Mourot et al.\ 2015; \toverify: DOI unresolved) with a functional outcome. The threshold between active recovery and additional connective tissue stress is not empirically defined for any climbing grip position.

\subsection*{Evidence Gaps and Pending Sections}

Across all domains, the critical shared gaps are: absence of female-cohort climbing data; lack of long-duration ($>$10 week) trials with concurrent injury tracking; and no trial measuring grade improvement as a primary outcome. Performance predictors (finger strength, pull-up metrics, and their correlation with climbing grade), equipment ROI, scientific measurement methodology (Tindeq Critical Force, W', force asymmetry), and skincare are addressed in forthcoming sections.

Until climbing-specific high-quality RCTs exist for these domains, practitioner folklore should be applied as biologically plausible guidance calibrated to individual tolerance, not as protocol with established dosing certainty.


\section*{Conflict of Interest}
The author declares no conflict of interest.

%---------------------------------------------------------------------
\appendix
\section{\toverify Claim Tracker}
\label{app:toverify}

The claims below carry the \toverify label. They must not be presented as settled guidance. Each entry requires DOI resolver confirmation and independent source verification before promotion to \evidencebacked, \bioplausible, or \folkloreverified.

\renewcommand{\arraystretch}{1.2}
\begin{longtable}{>{\raggedright}p{0.5cm}
                  >{\raggedright}p{4.0cm}
                  >{\raggedright}p{2.2cm}
                  >{\raggedright}p{2.2cm}
                  >{\raggedright}p{2.0cm}
                  >{\raggedright\arraybackslash}p{3.5cm}}
\caption{\toverify claim registry. Update when resolved.}\label{tab:toverify}\\
\toprule
\textbf{\#} & \textbf{Claim} & \textbf{Section} & \textbf{Current label} & \textbf{ESS} & \textbf{Required action} \\
\midrule
\endfirsthead
\toprule
\textbf{\#} & \textbf{Claim} & \textbf{Section} & \textbf{Current label} & \textbf{ESS} & \textbf{Required action} \\
\midrule
\endhead
\bottomrule
\endlastfoot

1 & Stien et al.\ (2023) volume/page conflict (40(3):875--889 vs.\ 40(1):179--191) & \S\ref{sec:weekly} & \toverify & --- & Confirm canonical journal volume/issue via PubMed PMID 37214263 \\
\addlinespace
2 & Hermans et al.\ (2022) single-source citation in weekly integration & \S\ref{sec:weekly} & \toverify & --- & Verify DOI 10.3389/fspor.2022.888158 independently; confirm $n$, protocol \\
\addlinespace
3 & Creatine body-mass penalty claim in climbing performance & \S\ref{sec:supplements} & \toverify & 1 & Identify $\geq$3 independent climbing-community sources beyond r/climbharder and supplement blogs \\
\addlinespace
4 & Omega-3 strengthens tendons / prevents pulley injuries & \S\ref{sec:supplements} & \toverify & 0 & Identify $\geq$3 independent climbing-community sources; current claim unsupported by controlled evidence \\
\addlinespace
5 & Magnesium prevents forearm pump in euvolaemic climbers & \S\ref{sec:supplements} & \toverify & 0 & Identify $\geq$3 independent climbing-community sources; verify specific JAMA cramp RCT DOI \\
\addlinespace
6 & Beta-alanine eliminates forearm pump & \S\ref{sec:supplements} & \toverify & 0 & Identify $\geq$3 independent climbing-community sources; Sas-Nowosielski 2021 shows partial improvement only \\
\addlinespace
7 & Return-to-climbing timeline after flappers (3--14\,d by severity) & \S\ref{sec:injury} & \toverify & 1 & Identify $\geq$3 independent named climbing sources with specific timelines \\
\addlinespace
8 & HRV readiness thresholds ($\pm$0.5 SD rule) for climbing & \S\ref{sec:recovery} & \toverify & 1 & Identify $\geq$3 independent sources; current derivation is from endurance coaching, not climbing \\
\addlinespace
9 & Active recovery formats (jug laps, ARC traversing) for climbers & \S\ref{sec:recovery} & \toverify & 1 & Identify $\geq$3 independent named sources beyond Anderson \& Anderson 2014 \\

\end{longtable}

\bibliographystyle{alpha}
\bibliography{bibliography}

\end{document}
